% Options for packages loaded elsewhere
% Options for packages loaded elsewhere
\PassOptionsToPackage{unicode}{hyperref}
\PassOptionsToPackage{hyphens}{url}
\PassOptionsToPackage{dvipsnames,svgnames,x11names}{xcolor}
%
\documentclass[
  letterpaper,
  DIV=11,
  numbers=noendperiod]{scrreprt}
\usepackage{xcolor}
\usepackage{amsmath,amssymb}
\setcounter{secnumdepth}{5}
\usepackage{iftex}
\ifPDFTeX
  \usepackage[T1]{fontenc}
  \usepackage[utf8]{inputenc}
  \usepackage{textcomp} % provide euro and other symbols
\else % if luatex or xetex
  \usepackage{unicode-math} % this also loads fontspec
  \defaultfontfeatures{Scale=MatchLowercase}
  \defaultfontfeatures[\rmfamily]{Ligatures=TeX,Scale=1}
\fi
\usepackage{lmodern}
\ifPDFTeX\else
  % xetex/luatex font selection
  \setmainfont[]{Source Serif 4}
  \setsansfont[]{Libre Franklin}
  \setmonofont[]{JetBrains Mono}
\fi
% Use upquote if available, for straight quotes in verbatim environments
\IfFileExists{upquote.sty}{\usepackage{upquote}}{}
\IfFileExists{microtype.sty}{% use microtype if available
  \usepackage[]{microtype}
  \UseMicrotypeSet[protrusion]{basicmath} % disable protrusion for tt fonts
}{}
\makeatletter
\@ifundefined{KOMAClassName}{% if non-KOMA class
  \IfFileExists{parskip.sty}{%
    \usepackage{parskip}
  }{% else
    \setlength{\parindent}{0pt}
    \setlength{\parskip}{6pt plus 2pt minus 1pt}}
}{% if KOMA class
  \KOMAoptions{parskip=half}}
\makeatother
% Make \paragraph and \subparagraph free-standing
\makeatletter
\ifx\paragraph\undefined\else
  \let\oldparagraph\paragraph
  \renewcommand{\paragraph}{
    \@ifstar
      \xxxParagraphStar
      \xxxParagraphNoStar
  }
  \newcommand{\xxxParagraphStar}[1]{\oldparagraph*{#1}\mbox{}}
  \newcommand{\xxxParagraphNoStar}[1]{\oldparagraph{#1}\mbox{}}
\fi
\ifx\subparagraph\undefined\else
  \let\oldsubparagraph\subparagraph
  \renewcommand{\subparagraph}{
    \@ifstar
      \xxxSubParagraphStar
      \xxxSubParagraphNoStar
  }
  \newcommand{\xxxSubParagraphStar}[1]{\oldsubparagraph*{#1}\mbox{}}
  \newcommand{\xxxSubParagraphNoStar}[1]{\oldsubparagraph{#1}\mbox{}}
\fi
\makeatother

\usepackage{color}
\usepackage{fancyvrb}
\newcommand{\VerbBar}{|}
\newcommand{\VERB}{\Verb[commandchars=\\\{\}]}
\DefineVerbatimEnvironment{Highlighting}{Verbatim}{commandchars=\\\{\}}
% Add ',fontsize=\small' for more characters per line
\usepackage{framed}
\definecolor{shadecolor}{RGB}{241,243,245}
\newenvironment{Shaded}{\begin{snugshade}}{\end{snugshade}}
\newcommand{\AlertTok}[1]{\textcolor[rgb]{0.68,0.00,0.00}{#1}}
\newcommand{\AnnotationTok}[1]{\textcolor[rgb]{0.37,0.37,0.37}{#1}}
\newcommand{\AttributeTok}[1]{\textcolor[rgb]{0.40,0.45,0.13}{#1}}
\newcommand{\BaseNTok}[1]{\textcolor[rgb]{0.68,0.00,0.00}{#1}}
\newcommand{\BuiltInTok}[1]{\textcolor[rgb]{0.00,0.23,0.31}{#1}}
\newcommand{\CharTok}[1]{\textcolor[rgb]{0.13,0.47,0.30}{#1}}
\newcommand{\CommentTok}[1]{\textcolor[rgb]{0.37,0.37,0.37}{#1}}
\newcommand{\CommentVarTok}[1]{\textcolor[rgb]{0.37,0.37,0.37}{\textit{#1}}}
\newcommand{\ConstantTok}[1]{\textcolor[rgb]{0.56,0.35,0.01}{#1}}
\newcommand{\ControlFlowTok}[1]{\textcolor[rgb]{0.00,0.23,0.31}{\textbf{#1}}}
\newcommand{\DataTypeTok}[1]{\textcolor[rgb]{0.68,0.00,0.00}{#1}}
\newcommand{\DecValTok}[1]{\textcolor[rgb]{0.68,0.00,0.00}{#1}}
\newcommand{\DocumentationTok}[1]{\textcolor[rgb]{0.37,0.37,0.37}{\textit{#1}}}
\newcommand{\ErrorTok}[1]{\textcolor[rgb]{0.68,0.00,0.00}{#1}}
\newcommand{\ExtensionTok}[1]{\textcolor[rgb]{0.00,0.23,0.31}{#1}}
\newcommand{\FloatTok}[1]{\textcolor[rgb]{0.68,0.00,0.00}{#1}}
\newcommand{\FunctionTok}[1]{\textcolor[rgb]{0.28,0.35,0.67}{#1}}
\newcommand{\ImportTok}[1]{\textcolor[rgb]{0.00,0.46,0.62}{#1}}
\newcommand{\InformationTok}[1]{\textcolor[rgb]{0.37,0.37,0.37}{#1}}
\newcommand{\KeywordTok}[1]{\textcolor[rgb]{0.00,0.23,0.31}{\textbf{#1}}}
\newcommand{\NormalTok}[1]{\textcolor[rgb]{0.00,0.23,0.31}{#1}}
\newcommand{\OperatorTok}[1]{\textcolor[rgb]{0.37,0.37,0.37}{#1}}
\newcommand{\OtherTok}[1]{\textcolor[rgb]{0.00,0.23,0.31}{#1}}
\newcommand{\PreprocessorTok}[1]{\textcolor[rgb]{0.68,0.00,0.00}{#1}}
\newcommand{\RegionMarkerTok}[1]{\textcolor[rgb]{0.00,0.23,0.31}{#1}}
\newcommand{\SpecialCharTok}[1]{\textcolor[rgb]{0.37,0.37,0.37}{#1}}
\newcommand{\SpecialStringTok}[1]{\textcolor[rgb]{0.13,0.47,0.30}{#1}}
\newcommand{\StringTok}[1]{\textcolor[rgb]{0.13,0.47,0.30}{#1}}
\newcommand{\VariableTok}[1]{\textcolor[rgb]{0.07,0.07,0.07}{#1}}
\newcommand{\VerbatimStringTok}[1]{\textcolor[rgb]{0.13,0.47,0.30}{#1}}
\newcommand{\WarningTok}[1]{\textcolor[rgb]{0.37,0.37,0.37}{\textit{#1}}}

\usepackage{longtable,booktabs,array}
\usepackage{calc} % for calculating minipage widths
% Correct order of tables after \paragraph or \subparagraph
\usepackage{etoolbox}
\makeatletter
\patchcmd\longtable{\par}{\if@noskipsec\mbox{}\fi\par}{}{}
\makeatother
% Allow footnotes in longtable head/foot
\IfFileExists{footnotehyper.sty}{\usepackage{footnotehyper}}{\usepackage{footnote}}
\makesavenoteenv{longtable}
\usepackage{graphicx}
\makeatletter
\newsavebox\pandoc@box
\newcommand*\pandocbounded[1]{% scales image to fit in text height/width
  \sbox\pandoc@box{#1}%
  \Gscale@div\@tempa{\textheight}{\dimexpr\ht\pandoc@box+\dp\pandoc@box\relax}%
  \Gscale@div\@tempb{\linewidth}{\wd\pandoc@box}%
  \ifdim\@tempb\p@<\@tempa\p@\let\@tempa\@tempb\fi% select the smaller of both
  \ifdim\@tempa\p@<\p@\scalebox{\@tempa}{\usebox\pandoc@box}%
  \else\usebox{\pandoc@box}%
  \fi%
}
% Set default figure placement to htbp
\def\fps@figure{htbp}
\makeatother


% definitions for citeproc citations
\NewDocumentCommand\citeproctext{}{}
\NewDocumentCommand\citeproc{mm}{%
  \begingroup\def\citeproctext{#2}\cite{#1}\endgroup}
\makeatletter
 % allow citations to break across lines
 \let\@cite@ofmt\@firstofone
 % avoid brackets around text for \cite:
 \def\@biblabel#1{}
 \def\@cite#1#2{{#1\if@tempswa , #2\fi}}
\makeatother
\newlength{\cslhangindent}
\setlength{\cslhangindent}{1.5em}
\newlength{\csllabelwidth}
\setlength{\csllabelwidth}{3em}
\newenvironment{CSLReferences}[2] % #1 hanging-indent, #2 entry-spacing
 {\begin{list}{}{%
  \setlength{\itemindent}{0pt}
  \setlength{\leftmargin}{0pt}
  \setlength{\parsep}{0pt}
  % turn on hanging indent if param 1 is 1
  \ifodd #1
   \setlength{\leftmargin}{\cslhangindent}
   \setlength{\itemindent}{-1\cslhangindent}
  \fi
  % set entry spacing
  \setlength{\itemsep}{#2\baselineskip}}}
 {\end{list}}
\usepackage{calc}
\newcommand{\CSLBlock}[1]{\hfill\break\parbox[t]{\linewidth}{\strut\ignorespaces#1\strut}}
\newcommand{\CSLLeftMargin}[1]{\parbox[t]{\csllabelwidth}{\strut#1\strut}}
\newcommand{\CSLRightInline}[1]{\parbox[t]{\linewidth - \csllabelwidth}{\strut#1\strut}}
\newcommand{\CSLIndent}[1]{\hspace{\cslhangindent}#1}



\setlength{\emergencystretch}{3em} % prevent overfull lines

\providecommand{\tightlist}{%
  \setlength{\itemsep}{0pt}\setlength{\parskip}{0pt}}



 


% Brand colors and spacing for the PDF build
\usepackage{xcolor}
\definecolor{ink}{HTML}{1B2A4A}
\definecolor{civicblue}{HTML}{2166AC}
\definecolor{teal}{HTML}{2A7F7F}
\addtokomafont{disposition}{\color{ink}}   % headings in ink
\addtokomafont{title}{\color{ink}}
\usepackage{microtype}
\usepackage{setspace}
\onehalfspacing
\KOMAoption{captions}{tableheading}
\makeatletter
\@ifpackageloaded{tcolorbox}{}{\usepackage[skins,breakable]{tcolorbox}}
\@ifpackageloaded{fontawesome5}{}{\usepackage{fontawesome5}}
\definecolor{quarto-callout-color}{HTML}{909090}
\definecolor{quarto-callout-note-color}{HTML}{0758E5}
\definecolor{quarto-callout-important-color}{HTML}{CC1914}
\definecolor{quarto-callout-warning-color}{HTML}{EB9113}
\definecolor{quarto-callout-tip-color}{HTML}{00A047}
\definecolor{quarto-callout-caution-color}{HTML}{FC5300}
\definecolor{quarto-callout-color-frame}{HTML}{acacac}
\definecolor{quarto-callout-note-color-frame}{HTML}{4582ec}
\definecolor{quarto-callout-important-color-frame}{HTML}{d9534f}
\definecolor{quarto-callout-warning-color-frame}{HTML}{f0ad4e}
\definecolor{quarto-callout-tip-color-frame}{HTML}{02b875}
\definecolor{quarto-callout-caution-color-frame}{HTML}{fd7e14}
\makeatother
\makeatletter
\@ifpackageloaded{bookmark}{}{\usepackage{bookmark}}
\makeatother
\makeatletter
\@ifpackageloaded{caption}{}{\usepackage{caption}}
\AtBeginDocument{%
\ifdefined\contentsname
  \renewcommand*\contentsname{Table of contents}
\else
  \newcommand\contentsname{Table of contents}
\fi
\ifdefined\listfigurename
  \renewcommand*\listfigurename{List of Figures}
\else
  \newcommand\listfigurename{List of Figures}
\fi
\ifdefined\listtablename
  \renewcommand*\listtablename{List of Tables}
\else
  \newcommand\listtablename{List of Tables}
\fi
\ifdefined\figurename
  \renewcommand*\figurename{Figure}
\else
  \newcommand\figurename{Figure}
\fi
\ifdefined\tablename
  \renewcommand*\tablename{Table}
\else
  \newcommand\tablename{Table}
\fi
}
\@ifpackageloaded{float}{}{\usepackage{float}}
\floatstyle{ruled}
\@ifundefined{c@chapter}{\newfloat{codelisting}{h}{lop}}{\newfloat{codelisting}{h}{lop}[chapter]}
\floatname{codelisting}{Listing}
\newcommand*\listoflistings{\listof{codelisting}{List of Listings}}
\makeatother
\makeatletter
\makeatother
\makeatletter
\@ifpackageloaded{caption}{}{\usepackage{caption}}
\@ifpackageloaded{subcaption}{}{\usepackage{subcaption}}
\makeatother
\usepackage{bookmark}
\IfFileExists{xurl.sty}{\usepackage{xurl}}{} % add URL line breaks if available
\urlstyle{same}
\hypersetup{
  pdftitle={Creating Local-Level Geographic Datasets},
  pdfauthor={Aaron Schroeder},
  colorlinks=true,
  linkcolor={RoyalBlue},
  filecolor={Maroon},
  citecolor={Blue},
  urlcolor={Blue},
  pdfcreator={LaTeX via pandoc}}


\title{Creating Local-Level Geographic Datasets}
\usepackage{etoolbox}
\makeatletter
\providecommand{\subtitle}[1]{% add subtitle to \maketitle
  \apptocmd{\@title}{\par {\large #1 \par}}{}{}
}
\makeatother
\subtitle{A practical, illustrated guide for sub-county policy analysis}
\author{Aaron Schroeder}
\date{2026-06-04}
\begin{document}
\maketitle

\renewcommand*\contentsname{Table of contents}
{
\hypersetup{linkcolor=}
\setcounter{tocdepth}{2}
\tableofcontents
}

\bookmarksetup{startatroot}

\chapter{The Sub-County Data Gap}\label{sec-gap}

Local officials make decisions at the neighborhood level, but the data
they rely on is usually reported at the county level. A county-wide
broadband adoption rate, a county-wide income figure, a county-wide
health outcome: these statistics obscure the variation that actually
matters for policy. A neighborhood with near-universal high-speed
service and an adjacent neighborhood with almost none look identical in
a county aggregate. The officials who need to target investment, design
interventions, or justify budget requests are working blind.

The units that drive real decisions are not counties. They are
neighborhood associations, service districts, planning corridors, school
catchment zones, and utility territories. These are the geographies
where residents organize, where services are delivered, and where
disparities become visible. County-level data cannot answer the
questions local officials are actually asking.

\section{Why sub-county data is
scarce}\label{why-sub-county-data-is-scarce}

The scarcity of reliable sub-county data has three overlapping causes.
First, small-area estimates carry large statistical uncertainty: when a
sample is divided into dozens of small geographic units, the counts per
unit shrink, and margins of error grow until estimates become unreliable
{[}1{]}. Second, statistical agencies and data providers apply
geographic suppression or aggregation to protect individual
confidentiality; the smaller the area, the more likely a cell will be
suppressed {[}2{]}. Third, there is no standardized collection
infrastructure at the sub-county level. Unlike counties or census
tracts, neighborhood and civic-association boundaries are locally
defined, vary across jurisdictions, and change over time, so no national
data program routinely publishes estimates for them {[}3{]}.

The result is a persistent data gap between the geographic scale at
which data is collected and reported and the geographic scale at which
local government operates. Analysts working inside local government (and
researchers trying to support them) must construct the sub-county
datasets they need rather than retrieve them from a shelf {[}4{]}.

\section{What this guide
demonstrates}\label{what-this-guide-demonstrates}

This guide works through a concrete example: estimating mean household
income and fixed-broadband download speeds for each of Arlington County,
Virginia's \textbf{62 civic associations}, which together cover
approximately \textbf{109,528 households}. The two source datasets
(Ookla fixed-broadband speed-test tiles and American Community Survey
block-group income and household counts) arrive on entirely different
geographic footprints that match neither each other nor the
civic-association boundaries. The guide shows, step by step, how to
align those footprints through areal-weighted interpolation and produce
estimates at the policy-relevant geography.

The worked example is implemented in \textbf{Python}. An \textbf{R}
implementation is forthcoming.

\begin{tcolorbox}[enhanced jigsaw, toptitle=1mm, opacityback=0, rightrule=.15mm, colback=white, bottomrule=.15mm, colframe=quarto-callout-note-color-frame, colbacktitle=quarto-callout-note-color!10!white, title=\textcolor{quarto-callout-note-color}{\faInfo}\hspace{0.5em}{Who this guide is for}, left=2mm, bottomtitle=1mm, breakable, titlerule=0mm, coltitle=black, toprule=.15mm, arc=.35mm, opacitybacktitle=0.6, leftrule=.75mm]

\textbf{Local government analysts} will find a reproducible workflow
they can adapt to their own jurisdiction, data sources, and policy
geographies. The method generalizes beyond broadband and income to any
pair of spatially misaligned datasets.

\textbf{Decision-makers and program managers} can read the prose and
figures without engaging the code. The key takeaways are in the chapter
narratives and maps. The technical sections are clearly marked.

Some background in geographic data is helpful but not required. The
guide explains each concept as it arises.

\end{tcolorbox}

\bookmarksetup{startatroot}

\chapter{The Problem of Misaligned Boundaries}\label{sec-maup}

Every geographic dataset comes packaged in a particular set of spatial
units: the shapes that define where each observation applies. Census
estimates arrive at the block-group or tract level. Speed-test tiles
come in a fixed web-mercator grid. Administrative records attach to
parcels, precincts, or service zones. The problem is that the boundaries
used to collect and publish one dataset almost never match those used
for another, and neither typically matches the policy geography where
decisions get made.

\section{The Modifiable Areal Unit
Problem}\label{the-modifiable-areal-unit-problem}

This mismatch has a formal name in spatial statistics: the
\textbf{Modifiable Areal Unit Problem}, or MAUP {[}5{]}. The MAUP has
two related components. The \emph{scale effect} describes how
statistical summaries (means, correlations, rates) change when the same
underlying data is aggregated to larger or smaller units. The
\emph{zonation effect} describes how the same summary statistics change
when you redraw boundaries at the same scale but in a different
configuration. Neither effect is a modeling error; both are inherent
properties of spatially aggregated data.

The practical consequence for local policy work is significant: the
relationship you observe between two variables can change substantially
depending solely on the geographic units you use to measure it {[}6{]}.
A correlation between income and broadband access that appears strong at
the block-group level may weaken, strengthen, or reverse at the
civic-association level, not because the underlying relationship
changed, but because the boundary choice changed. Analysts who do not
account for this can draw conclusions that are artifacts of the data's
packaging rather than reflections of reality.

To see why, consider a single variable measured across one small area
and summarised under two district schemes (Figure~\ref{fig-maup}). The
underlying values are identical in both cases; only the boundaries
differ. One scheme cuts across the gradient and reports a nearly uniform
figure of about 220 Mbps in every district, erasing the variation
entirely. The other aligns with the gradient and reports districts
ranging from 138 to 302 Mbps. Same data, opposite conclusions, with
nothing changed but where the lines were drawn.

\begin{figure}[H]

\centering{

\includegraphics[width=1\linewidth,height=\textheight,keepaspectratio]{figures/fig_maup.png}

}

\caption{\label{fig-maup}The Modifiable Areal Unit Problem in miniature.
A single east-to-west speed gradient (left) is summarised under two
district schemes covering the identical area. Boundaries A (centre) cut
across the gradient, so every district averages about 220 Mbps and the
variation disappears. Boundaries B (right) align with the gradient, so
district averages range from 138 to 302 Mbps. The choice of geography is
an analytical decision, not a clerical one.}

\end{figure}%

\section{Three misaligned
geographies}\label{three-misaligned-geographies}

This guide works with three source geographies, none of which aligns
with the others.

\textbf{Ookla speed-test tiles} are delivered on a zoom-level-16
web-mercator grid. Each tile is approximately 610 meters on a side: a
fixed, global grid with no relationship to any political or
administrative boundary. Tiles do not respect county lines, block-group
edges, or neighborhood borders.

\textbf{ACS block groups} are the Census Bureau's smallest standard
geography for which the five-year American Community Survey publishes
reliable estimates. Block groups within Arlington County range widely in
area and population density; their boundaries follow streets, water
features, and other physical markers that predate any local policy unit.

\textbf{Arlington civic associations} are community-defined
neighborhoods and the county's formal channels for resident input and
local governance. Their 62 boundaries were drawn through local political
and historical processes, not by any statistical agency, and they bear
no systematic relationship to the Census hierarchy.

The figure below illustrates the transformation this guide performs:
taking source data packaged in two incompatible grids and producing
estimates on the policy-relevant civic-association geography.

\begin{figure}[H]

\centering{

\includegraphics[width=1\linewidth,height=\textheight,keepaspectratio]{figures/fig_transformation_3panel.png}

}

\caption{\label{fig-transform}Transforming misaligned source data (Ookla
tiles, ACS block groups) onto the policy-relevant civic-association
geography.}

\end{figure}%

As Figure~\ref{fig-transform} shows, a single civic association may
overlap dozens of Ookla tiles and several block groups simultaneously.
No one-to-one correspondence exists between any pair of these three
geographies. The method described in subsequent chapters, areal-weighted
interpolation, handles this overlap systematically by distributing each
source observation across target units in proportion to the share of its
area that falls within each target unit.

The key assumption underlying this approach is that the quantity of
interest is distributed uniformly within each source unit. That
assumption is never perfectly true, but it is a principled and auditable
starting point, and far preferable to ignoring the boundary problem
altogether.

\bookmarksetup{startatroot}

\chapter{The Data}\label{sec-data}

This chapter describes the three datasets used in the worked example:
Ookla fixed-broadband speed-test tiles, American Community Survey
block-group income and household counts, and Arlington County
civic-association boundaries. Each is publicly available.

\section{Ookla Speed-Test Tiles}\label{ookla-speed-test-tiles}

Ookla publishes quarterly snapshots of aggregated speed-test performance
through its open-data program {[}7{]}, {[}8{]}. The fixed-broadband
dataset contains one record per zoom-level-16 web-mercator tile, a
global grid in which each cell is approximately 610 meters on a side at
mid-latitudes. For each tile that received at least one test in a
quarter, the dataset reports the median download speed (Mbps), median
upload speed (Mbps), and number of tests that contributed to those
medians.

For this guide, the relevant quarter is Q1 2021, chosen to align with
the ACS 2021 five-year estimates. Download speed and test count are the
primary variables of interest. Test count matters because tiles with
very few tests carry more uncertainty; the method chapter discusses how
this is handled.

Ookla speed-test data reflects the experience of users who actively ran
a test at Speedtest.net or through an embedded Speedtest SDK during the
quarter. This is not a random sample of all internet users, and coverage
varies with population density and smartphone penetration. These
limitations are discussed further in the limitations chapter.

\begin{figure}[H]

\centering{

\includegraphics[width=0.9\linewidth,height=\textheight,keepaspectratio]{figures/fig_ookla_tiles.png}

}

\caption{\label{fig-tiles}Ookla fixed-broadband speed-test tiles over
Arlington (download Mbps).}

\end{figure}%

\section{American Community Survey}\label{american-community-survey}

The American Community Survey is the Census Bureau's ongoing household
survey, published annually as one-year and five-year estimates. The
five-year estimates average data collected over a rolling 60-month
window and are the appropriate choice for small geographies where
single-year samples are too thin to support reliable estimates {[}9{]},
{[}10{]}.

This guide uses two block-group-level variables from the ACS 2021
five-year release:

\begin{itemize}
\tightlist
\item
  \textbf{Table B19025:} Aggregate household income. This is the sum of
  all household incomes within the block group, expressed in dollars.
  Unlike the median, aggregate income is an additive quantity: totals
  from sub-areas can be summed to produce a total for any containing
  area.
\item
  \textbf{Table B11001:} Household count. The number of occupied housing
  units in the block group.
\end{itemize}

We use these two counts together rather than median income because mean
income for any custom geography can be derived directly from them:
divide aggregate income by household count. This arithmetic is only
valid for additive counts, not medians. Using aggregate income and
household counts preserves the ability to produce correct estimates at
any aggregated geographic unit.

\section{Arlington Civic
Associations}\label{arlington-civic-associations}

Arlington County, Virginia is divided into 62 civic associations:
community-defined neighborhoods that are the county's primary formal
mechanism for organized resident input {[}11{]}, {[}12{]}. Each
association elects officers and engages with county government on land
use, transportation, infrastructure, and community services. Membership
in the Arlington County Civic Federation, which coordinates the
associations, is the standard means by which neighborhoods participate
in county decision-making.

The civic-association boundaries used here come from Arlington County's
GIS open data portal {[}13{]}. The 62 associations cover the entire
county and tile it without overlap, making them a complete and
non-overlapping partition of Arlington's approximately 109,528
households. This makes them an ideal policy geography: they are locally
meaningful, formally recognized, and together exhaustive.

\begin{figure}[H]

\centering{

\includegraphics[width=0.9\linewidth,height=\textheight,keepaspectratio]{figures/fig_locator_civic.png}

}

\caption{\label{fig-locator}Arlington's 62 civic associations, the
policy-relevant target geography.}

\end{figure}%

Because civic associations are not part of the Census geographic
hierarchy, no federal statistical agency publishes estimates at this
level. Producing those estimates is precisely what the remainder of this
guide demonstrates.

\bookmarksetup{startatroot}

\chapter{The Core Method: Areal Interpolation}\label{sec-method}

The fundamental challenge of local-level data work is geometric: you
have numbers attached to one set of polygons and questions about a
different set of polygons. Areal interpolation is the family of methods
that bridges that gap. The specific technique used throughout this
guide, area-weighted redistribution, is the most transparent and widely
applicable member of that family {[}14{]}, {[}15{]}.

\section{How area-weighted redistribution
works}\label{how-area-weighted-redistribution-works}

The governing idea is simple: if a source area's polygon overlaps 40\%
of a target area's polygon, then 40\% of that source area's count is
assigned to that target area. More precisely, for each pair of source
and target polygons that intersect, the method computes the fractional
overlap (the ratio of intersection area to source area) and multiplies
the source value by that fraction. When a target polygon overlaps
multiple source polygons, the contributions from each source are summed
to produce the target estimate.

Formally, the weight \(w_{ij}\) applied when transferring a value from
source polygon \(i\) to target polygon \(j\) is:

\[w_{ij} = \frac{\text{area}(i \cap j)}{\text{area}(i)}\]

The weights across all target polygons that share area with source
polygon \(i\) sum to 1.0, so the total value in the source layer is
conserved. This is the key constraint that distinguishes redistribution
from other forms of spatial estimation: values are allocated, not
created.

The method rests on one explicit assumption: the variable of interest is
distributed uniformly within each source polygon. Population is not
perfectly uniform within a Census block group, and download speed is not
perfectly uniform across a 610-meter Ookla tile. But the assumption is
transparent, auditable, and consistent, and it systematically
outperforms simply ignoring the boundary problem altogether {[}16{]}.

\section{Extensive vs.~intensive
measures}\label{sec-extensive-intensive}

This is the most important concept in the entire chapter. Understanding
it determines whether your redistributed estimates are correct or subtly
wrong.

Geographic measures divide into two fundamentally different types based
on how they relate to the size of the area they describe.

\textbf{Extensive measures} (also called additive or count-like
measures) \emph{scale with area}. If you split a block group in two,
each half has roughly half the total households, roughly half the
aggregate income, roughly half the test count. Extensive measures can be
legitimately redistributed using area weights because the splitting and
summing is arithmetically valid. Examples: total households, aggregate
household income, total speed-test count.

\textbf{Intensive measures} (also called rate-like or derived measures)
\emph{do not scale with area}. Mean income is the same whether you
measure it for the whole block group or for a one-block sliver of it. If
you split a block group and assign each half 40\% of its mean income,
you have made a category error: you are treating a rate as if it were a
count. Examples: mean income, mean download speed, population density,
any ratio.

The practical rule is therefore:

\begin{quote}
\textbf{Redistribute counts. Derive rates afterward.}
\end{quote}

To produce a valid mean income for a civic association, you redistribute
\emph{aggregate income} (extensive) and \emph{household count}
(extensive) separately to each civic association, then divide:
\(\text{mean income} = \text{agg. income} \div \text{households}\). You
never redistribute mean income directly.

The same logic applies to a statistic that analysts sometimes want but
cannot obtain this way: \textbf{the median}. Median household income is
fundamentally a positional statistic (it describes the midpoint of a
distribution) and it has no additive decomposition. There is no
arithmetic operation on two partial medians that recovers the median of
their union. If you areally redistribute a median, the result is
numerically meaningless. This guide uses aggregate income precisely
because medians cannot be handled this way.

Getting this distinction wrong is the single most common error in areal
interpolation. It produces estimates that look plausible but are quietly
incorrect, typically biased toward area-weighted means of the
\emph{rate} rather than the \emph{population-weighted} mean that the
underlying counts would produce. In areas with uneven population density
(Arlington has denser urban corridors alongside lower-density suburban
neighborhoods), this error can be substantial.

Figure~\ref{fig-dataflow} summarizes the full redistribution workflow:
extensive counts enter the redistribution, area weights are applied, and
intensive rates are computed only at the end from the redistributed
count totals.

\begin{figure}[H]

\centering{

\includegraphics[width=1\linewidth,height=\textheight,keepaspectratio]{figures/diagrams/dataflow.png}

}

\caption{\label{fig-dataflow}The redistribution method: extensive counts
move area-weighted; intensive rates are derived from them.}

\end{figure}%

\section{\texorpdfstring{The \texttt{sdc-redistribute}
implementation}{The sdc-redistribute implementation}}\label{the-sdc-redistribute-implementation}

In practice, the redistribution is performed by the
\texttt{sdc\_redistribute} Python package, which handles the geometric
intersection, weight calculation, and column-wise multiplication in a
single call. The source geometry and one or more target geometries are
supplied as GeoJSON or GeoDataFrame objects; the caller specifies which
columns are extensive counts. The function returns a target-geometry
GeoDataFrame with redistributed count columns that can be combined or
divided to produce any desired derived measure.

\begin{tcolorbox}[enhanced jigsaw, toptitle=1mm, opacityback=0, rightrule=.15mm, colback=white, bottomrule=.15mm, colframe=quarto-callout-note-color-frame, colbacktitle=quarto-callout-note-color!10!white, title=\textcolor{quarto-callout-note-color}{\faInfo}\hspace{0.5em}{Python: the redistribution call}, left=2mm, bottomtitle=1mm, breakable, titlerule=0mm, coltitle=black, toprule=.15mm, arc=.35mm, opacitybacktitle=0.6, leftrule=.75mm]

\begin{Shaded}
\begin{Highlighting}[]
\ImportTok{from}\NormalTok{ sdc\_redistribute }\ImportTok{import}\NormalTok{ redistribute\_direct}

\NormalTok{result }\OperatorTok{=}\NormalTok{ redistribute\_direct(}
\NormalTok{    source\_df}\OperatorTok{=}\NormalTok{counts\_long,            }\CommentTok{\# geoid, year, measure, value}
\NormalTok{    source\_geo}\OperatorTok{=}\StringTok{"block\_groups.geojson"}\NormalTok{,}
\NormalTok{    target\_geos}\OperatorTok{=}\NormalTok{\{}\StringTok{"civic\_association"}\NormalTok{: }\StringTok{"civic\_assoc.geojson"}\NormalTok{\},}
\NormalTok{    count\_cols}\OperatorTok{=}\NormalTok{[}\StringTok{"agg\_income"}\NormalTok{, }\StringTok{"households"}\NormalTok{],   }\CommentTok{\# EXTENSIVE counts}
\NormalTok{    source\_id}\OperatorTok{=}\StringTok{"geoid"}\NormalTok{,}
\NormalTok{)}
\CommentTok{\# mean income (INTENSIVE) is derived afterwards:}
\CommentTok{\#   mean\_income = agg\_income / households}
\end{Highlighting}
\end{Shaded}

\end{tcolorbox}

The collapsible snippet above shows the essential pattern: only
extensive counts are passed to \texttt{count\_cols}. The intensive
quantity, mean income, is not a parameter to the function; it is
computed by the analyst on the output. This design makes the
extensive-vs-intensive distinction architecturally explicit: the
function cannot produce a rate; it can only redistribute counts.

The worked example in Chapter~\ref{sec-example} applies this pattern
twice in sequence, once for income and once for broadband speed, and
then validates that the redistributed totals are internally consistent
before proceeding to analysis.

\bookmarksetup{startatroot}

\chapter{Worked Example: Arlington, Step by Step}\label{sec-example}

This chapter walks the full pipeline from raw source files to validated
civic-association estimates. The five stages (acquire, redistribute
income, redistribute broadband, combine, and validate) correspond
directly to the pipeline modules in the companion repository. By the
end, each of Arlington's 62 civic associations has an estimated mean
household income and an estimated mean download speed derived from the
same underlying geography and the same redistribution logic.

The complete runnable pipeline is a single command:
\texttt{uv\ run\ python\ -m\ pipeline.run}.

\section{Stage 1: Acquire the source
data}\label{stage-1-acquire-the-source-data}

Three datasets must be in hand before any redistribution can take place.

\textbf{Census block-group boundaries} are fetched programmatically
using \texttt{pygris}, a Python wrapper around the Census Bureau's
TIGER/Line shapefiles API. The call retrieves the block-group polygons
for Arlington County (FIPS 51013) at the 2021 vintage, which aligns with
the ACS five-year estimates used for income. The resulting GeoDataFrame
provides the source geometry for the ACS redistribution step.

\textbf{ACS income variables} are retrieved through the Census API using
a Census API key. The two variables used are B19025\_001E (aggregate
household income in the past 12 months) and B11001\_001E (number of
occupied housing units, i.e., households). Both are block-group-level
estimates from the ACS 2021 five-year release. Neither is a median; both
are extensive counts that can be legitimately redistributed. A Census
API key is required; the pipeline reads it from the environment variable
\texttt{CENSUS\_API\_KEY}.

\textbf{Ookla fixed-broadband speed-test tiles} are downloaded from
Ookla's public Amazon S3 bucket, which hosts quarterly snapshots of the
aggregated tile data as geoparquet files. The pipeline fetches the Q1
2021 tile for the relevant S3 path, clips it to the Arlington bounding
box, and retains two columns: \texttt{tests} (the count of speed tests
in each tile during the quarter) and \texttt{avg\_d\_kbps} (the average
download speed across those tests, in kilobits per second). Both are
used in the broadband redistribution below.

With these three inputs in hand (block-group polygons, ACS counts, and
Ookla tiles), the pipeline is ready to redistribute.

\section{Stage 2: Redistribute
income}\label{stage-2-redistribute-income}

Income follows the pattern described in
Section~\ref{sec-extensive-intensive}. The two ACS variables, aggregate
household income and household count, are extensive measures. Both are
redistributed independently to civic associations using area-weighted
interpolation. The redistribution preserves their totals: whatever
aggregate income exists across all block groups flows, in its entirety,
to the set of civic associations (since civic associations tile the
county without gaps or overlaps).

Figure~\ref{fig-acs} illustrates the transformation. The source (block
groups) carries two count columns, and the target (civic associations)
receives redistributed versions of both.

\begin{figure}[H]

\centering{

\includegraphics[width=0.95\linewidth,height=\textheight,keepaspectratio]{figures/diagrams/acs_transform.png}

}

\caption{\label{fig-acs}Income transform: redistribute aggregate income
+ households, then derive mean income.}

\end{figure}%

Once redistribution is complete, mean household income for each civic
association is a simple arithmetic step:

\[\text{mean income}_j = \frac{\text{agg. income}_j}{\text{households}_j}\]

The result is a population-weighted mean: larger block groups (more
households) contribute more to the civic association's estimate than
smaller ones, scaled continuously by the fraction of their area that
falls within the association's boundary.

\section{Stage 3: Redistribute
broadband}\label{stage-3-redistribute-broadband}

Broadband speed requires one additional preparatory step because
download speed, like income, is an intensive measure. The variable
reported in the Ookla tiles, \texttt{avg\_d\_kbps}, cannot be directly
redistributed. The count that makes it meaningful, \texttt{tests}, can.

The solution is to reconstruct the extensive quantity that, when divided
by test count, recovers average speed. For each tile:

\[\text{speed} \times \text{tests} = \text{total download kilobits (proxy)}\]

This product is an extensive count: if a tile's area is split in two,
each half contributed roughly half the total download kilobits. It can
be redistributed. The test count is also extensive and can be
redistributed. After both are moved to civic associations, the intensive
rate is derived:

\[\text{download speed}_j = \frac{(\text{speed} \times \text{tests})_j}{\text{tests}_j} \div 1000\]

The division by 1,000 converts kilobits per second to megabits per
second for presentation.

Figure~\ref{fig-ookla} shows the broadband transform. Notice the
structural parallel with the income transform in Figure~\ref{fig-acs}:
in both cases, the intensive output (mean income, mean speed) appears
only after redistribution, never before.

\begin{figure}[H]

\centering{

\includegraphics[width=0.95\linewidth,height=\textheight,keepaspectratio]{figures/diagrams/ookla_transform.png}

}

\caption{\label{fig-ookla}Broadband transform: redistribute tests and
speed×tests, then derive count-weighted speed.}

\end{figure}%

\begin{tcolorbox}[enhanced jigsaw, toptitle=1mm, opacityback=0, rightrule=.15mm, colback=white, bottomrule=.15mm, colframe=quarto-callout-note-color-frame, colbacktitle=quarto-callout-note-color!10!white, title=\textcolor{quarto-callout-note-color}{\faInfo}\hspace{0.5em}{Python: count-weighting broadband speed}, left=2mm, bottomtitle=1mm, breakable, titlerule=0mm, coltitle=black, toprule=.15mm, arc=.35mm, opacitybacktitle=0.6, leftrule=.75mm]

\begin{Shaded}
\begin{Highlighting}[]
\CommentTok{\# Speed is INTENSIVE: redistribute counts, then derive the rate.}
\NormalTok{tiles[}\StringTok{"d\_product"}\NormalTok{] }\OperatorTok{=}\NormalTok{ tiles[}\StringTok{"avg\_d\_kbps"}\NormalTok{] }\OperatorTok{*}\NormalTok{ tiles[}\StringTok{"tests"}\NormalTok{]   }\CommentTok{\# extensive}
\CommentTok{\# ... redistribute tests + d\_product to civic associations ...}
\NormalTok{civic[}\StringTok{"download\_mbps"}\NormalTok{] }\OperatorTok{=}\NormalTok{ (civic[}\StringTok{"d\_product"}\NormalTok{] }\OperatorTok{/}\NormalTok{ civic[}\StringTok{"tests"}\NormalTok{]) }\OperatorTok{/} \DecValTok{1000}
\end{Highlighting}
\end{Shaded}

\end{tcolorbox}

The collapsible snippet above captures the key lines. Creating
\texttt{d\_product} before redistribution is the move that makes
everything else valid. An analyst who skipped this step and
redistributed \texttt{avg\_d\_kbps} directly would produce a speed
estimate weighted by \emph{area}, not by \emph{test count}: a different
and less meaningful quantity that systematically overstates speed in
low-density tiles with few tests.

\section{Stage 4: Combine}\label{stage-4-combine}

With both redistributions complete, the income and broadband results are
joined to the civic-association GeoDataFrame on the association
identifier. The four intermediate columns (redistributed aggregate
income, redistributed households, redistributed test count, and
redistributed speed-product) are used to compute the two final output
columns: \texttt{mean\_income} and \texttt{download\_mbps}. All 62 civic
associations receive a value for both measures.

\section{Stage 5: Validate}\label{stage-5-validate}

Before proceeding to analysis or visualization, the pipeline verifies
one internal consistency check: household totals should be preserved by
the redistribution. The sum of redistributed households across all 62
civic associations is compared to the sum of ACS household counts across
all block groups. In a well-functioning redistribution over a complete,
non-overlapping partition of the county (which Arlington's civic
associations are), these totals should be identical up to floating-point
rounding.

In practice, the redistributed household total is preserved within
\textbf{2\%} of the source total, a result of minor edge effects at the
county boundary where tile clipping introduces small geometric
imprecision. This tolerance is well within the margin of the ACS
sampling error itself, and it confirms that the redistribution has not
created or destroyed households at any meaningful scale.

Analysts applying this method to their own geographies should perform
this same check. A discrepancy larger than a few percent typically
indicates either a non-exhaustive target partition (gaps between target
polygons), a coordinate-reference-system mismatch between source and
target layers, or a geometric validity problem in one of the input
files.

With validation passed, the pipeline writes its output and the data is
ready for mapping and analysis, the subject of the next chapter.

\bookmarksetup{startatroot}

\chapter{Results \& Interpretation}\label{sec-results}

The pipeline described in the preceding chapters produces two estimates
for each of Arlington's 62 civic associations: a mean household income
derived from ACS block-group counts, and a mean download speed derived
from Ookla speed-test tiles. This chapter presents both distributions,
compares them, and draws out the central finding, which runs counter to
the most common assumption analysts bring to digital-divide work.

\section{Income distribution}\label{income-distribution}

Mean household income across the 62 associations spans from roughly
\textbf{\$73,000} (Arlington Mill) to \textbf{\$468,000} (Gulf Branch),
with a mean of association means near \textbf{\$215,600} and a median
near \textbf{\$191,500}. The roughly \textbf{109,500} households
represented in the data are not evenly spread. Arlington's geography
concentrates its highest incomes in low-density, single-family
neighbourhoods in the northern and western portions of the county: Gulf
Branch, Rivercrest, Bellevue Forest, and Old Glebe all exceed \$447,000.
Denser, more mixed-use corridors in the south and east (Columbia
Heights, Arlington Mill, Forest Glen) cluster toward the lower end.

\begin{figure}[H]

\centering{

\includegraphics[width=0.85\linewidth,height=\textheight,keepaspectratio]{figures/map_income.png}

}

\caption{\label{fig-income}Mean household income by civic association.}

\end{figure}%

The spatial pattern in Figure~\ref{fig-income} reflects Arlington's
well-documented socio-spatial structure: income rises sharply as one
moves away from the Rosslyn--Ballston and Route 1 corridors into the
low-density residential tracts that border Fairfax and Falls Church
{[}17{]}. The spread is wide for a jurisdiction that ranks among the
wealthiest counties in the United States, which makes the county an
instructive case: if income structured broadband access anywhere, it
should be detectable here.

\section{Broadband speed
distribution}\label{broadband-speed-distribution}

Mean download speed across the same 62 associations ranges from
\textbf{100.8 Mbps} (Foxcroft Heights) to \textbf{320.8 Mbps} (Columbia
Forest), with a county-wide mean of approximately \textbf{246.8 Mbps}
and a median of \textbf{251.9 Mbps}. Mean upload speed averages roughly
\textbf{105 Mbps}. Even at the low end, every association exceeds the
FCC's legacy 25/3 Mbps benchmark; the real differentiation occurs at the
higher tiers that support heavy simultaneous use.

\begin{figure}[H]

\centering{

\includegraphics[width=0.85\linewidth,height=\textheight,keepaspectratio]{figures/map_speed.png}

}

\caption{\label{fig-speed}Broadband download speed by civic
association.}

\end{figure}%

The spatial pattern in Figure~\ref{fig-speed} is striking and largely
inverts the income map. The fastest associations are concentrated in the
denser southern and central corridors. Columbia Forest reaches 320.8
Mbps; Long Branch Creek 307.9 Mbps; \textbf{Arlington Mill, the
lowest-income association in the dataset at approximately \$73,000,
records the third-fastest speed at 307.0 Mbps}. Fairlington follows at
293.4 Mbps. Conversely, the slowest associations (Foxcroft Heights at
100.8 Mbps, Arlingwood at 117.3 Mbps, Dominion Hills at 176.2 Mbps) are
not low-income enclaves; they are mid-to-high-income, low-density,
single-family neighbourhoods in the county's outer northern tier.

\section{The central finding: income and speed are
decoupled}\label{the-central-finding-income-and-speed-are-decoupled}

The Pearson correlation between mean household income and mean download
speed across the 62 civic associations is \textbf{r = −0.11}. This value
is substantively indistinguishable from zero. Income and broadband speed
are essentially \textbf{decoupled} in Arlington.

\begin{figure}[H]

\centering{

\includegraphics[width=0.85\linewidth,height=\textheight,keepaspectratio]{figures/scatter_income_speed.png}

}

\caption{\label{fig-scatter}Income vs.~broadband download speed across
the 62 civic associations (r = −0.11).}

\end{figure}%

The scatter plot in Figure~\ref{fig-scatter} makes the absence of a
relationship visually plain. There is no line, positive or negative,
that summarises the cloud in any useful way. Analysts who arrive at this
data expecting the standard ``rich neighbourhoods get faster internet''
finding will not find it here. The naive expectation, plausible on its
face and documented in national-scale studies of access gaps {[}18{]},
does not hold inside a uniformly affluent jurisdiction where income
variation is high but where the key driver of infrastructure investment
operates on a different dimension entirely.

That dimension is \textbf{housing density}. Dense multifamily corridors
generate the subscriber concentration that makes competitive fibre and
cable deployment commercially attractive to internet service providers.
Arlington's Rosslyn--Ballston spine and the Route 1 corridor have both
density and the infrastructure investment that follows from it.
Low-density single-family enclaves in northern Arlington have neither
the building stock to attract the same level of competitive investment
nor the resulting speeds, regardless of how wealthy their residents are.

\section{Income-to-speed ratio}\label{income-to-speed-ratio}

To quantify the mismatch between income and infrastructure, consider the
income-to-speed ratio: dollars of mean household income per Mbps of mean
download speed. Across the 62 associations this ratio ranges from
\textbf{\$238 per Mbps} (Arlington Mill, high speed, lower income) to
\textbf{\$2,490 per Mbps} (Arlingwood, low speed, higher income), more
than a ten-fold span.

\begin{figure}[H]

\centering{

\includegraphics[width=0.85\linewidth,height=\textheight,keepaspectratio]{figures/map_ratio.png}

}

\caption{\label{fig-ratio}Income-to-speed ratio (dollars of household
income per Mbps of download speed).}

\end{figure}%

The map in Figure~\ref{fig-ratio} is effectively a map of which
associations are getting the least broadband for their economic
standing. The highest-ratio associations (Arlingwood at \$2,490,
Bellevue Forest at \$2,107, Gulf Branch at \$1,752, Old Glebe at
\$1,713) are all wealthy, low-density, northern-tier communities. They
are not suffering a crisis of unaffordability; they are simply not
served by the same density of competitive infrastructure as their
southern counterparts. The lowest-ratio association, Arlington Mill, has
the strongest infrastructure relative to its income level precisely
because dense multifamily housing attracted competitive carriers.

\section{Bivariate synthesis}\label{bivariate-synthesis}

Classifying each association by income tercile (low/mid/high) and speed
tercile (low/mid/high) produces a three-by-three contingency table that
summarises the joint distribution. The cell counts are:

\begin{longtable}[]{@{}llll@{}}
\toprule\noalign{}
& Low speed & Mid speed & High speed \\
\midrule\noalign{}
\endhead
\bottomrule\noalign{}
\endlastfoot
\textbf{Low income} & 4 & 8 & 9 \\
\textbf{Mid income} & 6 & 7 & 7 \\
\textbf{High income} & 11 & 5 & 5 \\
\end{longtable}

\begin{figure}[H]

\centering{

\includegraphics[width=0.85\linewidth,height=\textheight,keepaspectratio]{figures/map_bivariate.png}

}

\caption{\label{fig-bivariate}Bivariate map: income tercile × speed
tercile across the 62 civic associations.}

\end{figure}%

The \textbf{largest single cell is high-income / low-speed, with 11
associations}: Williamsburg, Maywood, Dominion Hills, Arlingwood,
Bellevue Forest, Donaldson Run, and their neighbours. This cell alone
falsifies the ``rich areas get fast internet'' hypothesis. At the
opposite corner, \textbf{9 associations fall in the low-income /
high-speed cell}: Columbia Heights, Forest Glen, Columbia Forest,
Westover Village, Long Branch Creek, and Arlington Mill, the dense
multifamily corridors already identified as speed leaders.

The bivariate map in Figure~\ref{fig-bivariate} makes the geography
legible. The high-income/low-speed cluster occupies the low-density
northern and western portions of the county. The low-income/high-speed
cluster wraps around the denser central and southern corridors. The
pattern is structurally coherent and consistent across all three ways of
visualising it: the scatter plot, the ratio map, and the bivariate class
map.

\section{Policy implications}\label{policy-implications}

The finding that \textbf{the digital-divide axis in Arlington is housing
density and infrastructure build-out, not household income}, has direct
implications for how jurisdictions should design broadband equity
interventions {[}19{]}, {[}20{]}.

Income-based subsidy targeting, the most common policy tool, would
direct resources toward low-income households while leaving the actual
speed gaps unaddressed. The households that experience the weakest
broadband performance in Arlington are, on average, high-income
households in low-density neighbourhoods. An income-tested programme
would systematically miss them.

More effective interventions would target the mechanism driving the gap:
the commercial calculus that makes dense corridors attractive to
competitive carriers but leaves single-family enclaves with fewer
options. This suggests policy levers such as right-of-way facilitation
for fibre providers in low-density zones, municipal broadband
franchising conditioned on area coverage, and infrastructure
requirements built into residential subdivision permitting. None of
these is calibrated to household income; all are calibrated to the
density-and-competition gap that the data actually reveal.

The Arlington case cannot be assumed to generalise directly to
jurisdictions with lower overall incomes or with genuine gaps in
physical access infrastructure. But the methodological lesson applies
everywhere: verify the relationship between income and speed empirically
before designing income-based interventions.

\bookmarksetup{startatroot}

\chapter{A Second Approach: Parcel-Based
Redistribution}\label{sec-parcels}

The area-weighted method introduced in Chapter~\ref{sec-method} works
well, but it rests on a simplifying assumption worth examining: that
each variable of interest is distributed uniformly across a block
group's entire polygon. In practice, a block group's polygon includes
roads, parks, parking lots, commercial strips, and other land that
contains no housing at all. For measures tied to where people actually
live, this matters. Dasymetric methods replace the uniform-distribution
assumption by using an ancillary layer (in this case, parcel locations)
to concentrate redistributed values where housing actually sits
{[}16{]}.

This chapter implements that alternative and then asks directly: how
much does the choice of method change the result?

\section{The parcel data}\label{the-parcel-data}

Arlington County's Master Housing Unit Database (MHUD) provides a
parcel-level inventory of every residential unit in the county. The
layer used here contains \textbf{34,340 parcels} carrying roughly
\textbf{127,000 housing units} in total. The vast majority of those
parcels (\textbf{33,665}) are single-family detached or attached
(SFD/SFA) structures, typically carrying one or two units apiece. The
remaining \textbf{675} parcels are classified as multifamily (MULTI),
and they account for a disproportionate share of the county's housing
stock: the multifamily parcels are concentrated along the
Rosslyn--Ballston corridor and the Route 1 spine, where high-rise and
mid-rise buildings each contribute hundreds of units to the total.

\begin{figure}[H]

\centering{

\includegraphics[width=0.9\linewidth,height=\textheight,keepaspectratio]{figures/map_parcels.png}

}

\caption{\label{fig-parcels}Arlington housing units by parcel. Point
size is unit count; the few multifamily parcels carry most of the
county's housing units, concentrated in the Rosslyn--Ballston and Route
1 corridors.}

\end{figure}%

Figure~\ref{fig-parcels} makes the spatial concentration vivid. A small
number of large dots, each representing a multifamily building, account
for a large fraction of the county's residential footprint. The
single-family parcels are numerous but individually small, spread across
the county's lower-density northern and western neighbourhoods. This
heterogeneity is precisely what dasymetric redistribution is designed to
exploit: instead of spreading a block group's household counts uniformly
across its entire polygon, we concentrate them where the housing units
are.

\section{The method: unit-weighted
redistribution}\label{the-method-unit-weighted-redistribution}

The implementation follows a straightforward unit-weighting scheme. Each
parcel is ``exploded'' into a number of identical point records equal to
its \texttt{Total\_Units} count: a 300-unit multifamily building
contributes 300 identical centroid points, while a single-family parcel
contributes one. That pool of unit points is then passed to
\texttt{redistribute\_parcels} from the \texttt{sdc-redistribute}
package, which uses the point locations as the dasymetric control layer.
The block-group aggregate income and household counts are redistributed
to civic associations in proportion to the number of unit points that
fall within each association, weighted by their source block-group
membership.

\begin{tcolorbox}[enhanced jigsaw, toptitle=1mm, opacityback=0, rightrule=.15mm, colback=white, bottomrule=.15mm, colframe=quarto-callout-note-color-frame, colbacktitle=quarto-callout-note-color!10!white, title=\textcolor{quarto-callout-note-color}{\faInfo}\hspace{0.5em}{Python: parcel redistribution}, left=2mm, bottomtitle=1mm, breakable, titlerule=0mm, coltitle=black, toprule=.15mm, arc=.35mm, opacitybacktitle=0.6, leftrule=.75mm]

\begin{Shaded}
\begin{Highlighting}[]
\ImportTok{from}\NormalTok{ sdc\_redistribute }\ImportTok{import}\NormalTok{ redistribute\_parcels}
\CommentTok{\# explode each parcel into Total\_Units identical points first, then:}
\NormalTok{result }\OperatorTok{=}\NormalTok{ redistribute\_parcels(}
\NormalTok{    source\_df}\OperatorTok{=}\NormalTok{counts\_long, parcel\_centroids}\OperatorTok{=}\NormalTok{unit\_points,}
\NormalTok{    source\_geo}\OperatorTok{=}\StringTok{"block\_groups.geojson"}\NormalTok{,}
\NormalTok{    target\_geos}\OperatorTok{=}\NormalTok{\{}\StringTok{"civic\_association"}\NormalTok{: }\StringTok{"civic\_assoc.geojson"}\NormalTok{\},}
\NormalTok{    count\_cols}\OperatorTok{=}\NormalTok{[}\StringTok{"agg\_income"}\NormalTok{, }\StringTok{"households"}\NormalTok{], source\_id}\OperatorTok{=}\StringTok{"geoid"}\NormalTok{,}
\NormalTok{)}
\end{Highlighting}
\end{Shaded}

\end{tcolorbox}

The logic is the same as the area-weighted approach (redistribute
extensive counts, derive intensive rates afterward), with one structural
difference: the implicit population surface has changed from a uniform
density within each block group to a density proportional to the spatial
distribution of housing units within that block group.

\section{Results: a close agreement with the area-weighted
estimates}\label{results-a-close-agreement-with-the-area-weighted-estimates}

\begin{figure}[H]

\centering{

\includegraphics[width=0.85\linewidth,height=\textheight,keepaspectratio]{figures/map_income_parcels.png}

}

\caption{\label{fig-income-parcels}Mean household income by civic
association, parcel-based redistribution.}

\end{figure}%

\begin{figure}[H]

\centering{

\includegraphics[width=0.85\linewidth,height=\textheight,keepaspectratio]{figures/map_income_diff.png}

}

\caption{\label{fig-income-diff}Difference in mean household income:
parcel minus area-weighted (dollars).}

\end{figure}%

Figure~\ref{fig-income-parcels} and Figure~\ref{fig-income-diff} tell a
consistent story. The parcel map looks nearly identical to its
area-weighted counterpart, and the difference map is dominated by values
close to zero.

\begin{figure}[H]

\centering{

\includegraphics[width=0.85\linewidth,height=\textheight,keepaspectratio]{figures/scatter_area_vs_parcel.png}

}

\caption{\label{fig-area-parcel}Area-weighted vs.~parcel-weighted mean
income across the 62 civic associations.}

\end{figure}%

The scatter plot in Figure~\ref{fig-area-parcel} makes the relationship
explicit: the two methods produce a Pearson r = 1.00 across the 62
associations. The largest per-association shift in either direction is
approximately ±\$22,000, a few percent on six-figure incomes. The three
biggest gainers under parcel weighting are \textbf{Arlington Forest
(+\$22,306)}, \textbf{Cherrydale (+\$15,347)}, and \textbf{Foxcroft
Heights (+\$13,047)}. The three biggest losers are \textbf{Lyon Park
(−\$6,118)}, \textbf{John M Langston (−\$4,960)}, and \textbf{Douglas
Park (−\$3,781)}.

Why do the methods agree so closely? Mean household income is a smooth,
intensive measure: it describes a characteristic of households, not a
count of things that accumulate over area. When a block group's income
is fairly evenly distributed among its housing stock, as it generally is
in Arlington at the block-group scale, then \emph{where} within the
block group the housing sits barely changes the weighted average that
reaches each civic association. Parcel placement shifts modest fractions
of the income count from one association to an adjacent one, but because
the income levels on both sides of any given boundary are similar, the
effect on the mean nearly cancels. The ±\$22,000 maximum shift
represents the cases where a multifamily building with an unusually
different household composition sits near a boundary: a real effect, but
a small one relative to the income range in the data.

\section{Leakage and the case for parcel weighting on count
measures}\label{leakage-and-the-case-for-parcel-weighting-on-count-measures}

One honest limitation deserves direct statement. Approximately
\textbf{3,032 households} (about \textbf{2.8\%} of the county total) sit
in parcels whose centroids fall outside any civic association boundary.
These households contribute to the block-group totals but are not
allocated to any association under the parcel method. The parcel
approach therefore slightly underestimates total households compared
with the area-weighted method, which allocates every square metre of a
block group to some association. Analysts using this method for official
estimates should be aware of this leakage and decide whether to
redistribute orphaned parcels to the nearest association or treat them
as a separate ``unassigned'' category.

For mean income, which is what this guide measures, the 2.8\% household
leakage has a negligible effect on the ratio. For \textbf{count
measures} (total population, total housing units), the answer is
different. Dasymetric and parcel-based redistribution earn their keep
precisely when the placement of housing within a block group drives the
allocation of counts across target boundaries. If you want to know how
many housing units fall within a specific civic association's polygon,
point-in-polygon counting on the parcel layer is more accurate than area
weighting; the difference can be substantial in associations with
compact boundaries or with multifamily buildings near their edges.

The natural next step in this direction is \textbf{household-size
weighting}: single-family and multifamily units have systematically
different average household sizes, and a redistribution that treats all
units as interchangeable conflates these. The \texttt{sdc-redistribute}
package's \texttt{redistribute\_parcel\_pums\_adj} function implements
this extension by drawing on Public Use Microdata Sample (PUMS)
household-size distributions, stratified by unit type, to weight
single-family parcels differently from multifamily ones. Applying that
method to Arlington's MHUD data, and assessing whether the
household-size correction changes the civic-association estimates
meaningfully, is a planned extension of this work.

For the income analysis in this guide, the two methods are
interchangeable in practice. The parcel approach is presented here to
establish the infrastructure and to make a broader point: method choice
matters most when the variable in question is a count and when housing
is unevenly distributed within source units. When the variable is a
smooth intensive measure like mean income, and when the redistribution
is operating at the block-group-to-neighbourhood scale, area weighting
and parcel weighting converge.

\bookmarksetup{startatroot}

\chapter{Limitations \& Uncertainty}\label{sec-limits}

The results in Chapter~\ref{sec-results} are derived estimates, not
direct measurements. Every step of the pipeline, from the source data to
the redistribution to the final ratio calculations, introduces
uncertainty. This chapter itemises the four most important sources and
explains what each implies for interpreting and extending the analysis.

\section{Area-weighting assumes uniform internal
distribution}\label{area-weighting-assumes-uniform-internal-distribution}

The core assumption of area-weighted interpolation is that the quantity
being redistributed is distributed uniformly within each source unit.
For ACS income, the source units are census block groups; for Ookla
broadband, the source units are 610-metre speed-test tiles. Wherever
that uniformity assumption is wrong, wherever households or speed-test
activity are spatially clustered within a unit, the redistribution will
misallocate a fraction of the value.

In practice, both income and population density vary within block
groups, particularly at the edges where block groups straddle the
civic-association boundary. The effect is largest for target
associations that are small relative to the source units they overlap,
and smallest for large associations that wholly contain many source
units. Analysts should treat estimates for the smallest associations
with additional caution and, where possible, validate against
parcel-level or point-level data.

This limitation is not unique to this workflow; it is inherent to any
area-weighted interpolation approach {[}14{]}, {[}21{]}. The
\texttt{redistribute\_direct} function used here, like all
area-weighting methods, cannot recover sub-unit spatial variation that
is not encoded in the source data.

\section{Ookla data reflects active test-takers, not all
households}\label{ookla-data-reflects-active-test-takers-not-all-households}

The Ookla speed-test tiles record the average performance of users who
actively ran a speed test during the measurement quarter. This
introduces a form of self-selection bias: test-takers are not a random
sample of all internet subscribers. People who run speed tests are more
likely to be technically engaged, more likely to be troubleshooting a
slow connection, and possibly more concentrated in particular building
types or demographic groups {[}22{]}.

The direction of this bias is not straightforwardly predictable. If
frustrated users test more frequently, the Ookla data may understate
typical speeds. If technically sophisticated users dominate the test
pool, the data may overstate typical speeds. For the purposes of a
comparative analysis across civic associations within a single county,
what matters is whether the bias is spatially systematic: whether the
test-taker population differs in structure across associations in ways
that would distort the between-association comparisons. That cannot be
fully evaluated without ground-truth data.

Users of this analysis should treat the download speeds as indicative of
infrastructure capacity at the high end of the user experience
distribution, not as estimates of median or typical household
experience.

\section{ACS income carries sampling
error}\label{acs-income-carries-sampling-error}

The ACS five-year block-group estimates are not administrative records;
they are survey estimates with associated margins of error. At the
block-group level, particularly for smaller block groups, the
coefficient of variation on income variables can be substantial. The
guide uses aggregate household income (B19025\_001E) and household count
(B11001\_001E) rather than median income precisely because both are
extensive, summable measures that survive redistribution with
interpretable units.

However, the mean income ultimately derived for each civic association
is still the ratio of two redistributed survey estimates. Both the
numerator (aggregate income) and denominator (household count) carry ACS
sampling uncertainty. In well-populated block groups this uncertainty is
modest; in smaller block groups it can be large enough to materially
affect the redistributed estimate. Arlington's relatively high response
rates and the use of five-year pooled estimates (2017--2021) mitigate
but do not eliminate this concern {[}23{]}.

The redistribution does not propagate formal uncertainty bounds. An
analyst who needs confidence intervals on the civic-association
estimates would need to request the block-group-level margins of error
from the Census API and carry them through the redistribution algebra, a
non-trivial extension that is beyond the scope of the current pipeline.

\section{\texorpdfstring{\texttt{redistribute\_direct} rescales measures
independently}{redistribute\_direct rescales measures independently}}\label{redistribute_direct-rescales-measures-independently}

The \texttt{redistribute\_direct} function rescales each source column
independently to ensure its total is preserved in the target geography.
When two columns are redistributed independently and then combined into
a ratio, as happens with both income (aggregate income / households) and
broadband (speed-product / tests), the rescaling applied to each column
need not be identical. Near the county boundary, where tile clipping and
block-group edge effects are most pronounced, the two columns may be
rescaled by slightly different factors, producing a ratio that is not
perfectly equivalent to what a joint redistribution would yield.

In the Arlington validation, the household total is preserved within
approximately 2\% of the source total, well within the ACS sampling
error and acceptable for a comparative county-level analysis. For
applications where ratio precision at the county boundary is critical,
analysts should consider masking associations that have more than a
threshold fraction of their area outside the source coverage, or running
a sensitivity check that compares estimates from independent and joint
redistribution paths.

\bookmarksetup{startatroot}

\chapter{Apply It to Your Own Jurisdiction}\label{sec-apply}

The Arlington analysis is a worked illustration of a general workflow.
Any jurisdiction with civic associations, planning districts, school
attendance zones, or other custom geographies can follow the same steps
to estimate income and broadband speed (or any other combination of an
ACS count variable and an Ookla speed-test variable) at the sub-county
level. This chapter distils the pipeline into a numbered recipe and
points to the tools needed to run it.

\section{The six-step recipe}\label{the-six-step-recipe}

\textbf{1. Define your target geography as a GeoJSON with a stable
identifier.} Your target polygons, the units you want estimates for,
must be stored as a GeoJSON (or any format readable by GeoPandas) and
must include a column that uniquely identifies each polygon. The
companion repository uses \texttt{geoid} by convention, but any string
or integer key works as long as it is stable and unique. The polygons
must be in a standard coordinate reference system (WGS 84 / EPSG:4326 is
the safe default) and should tile the jurisdiction without meaningful
gaps or overlaps. Gaps will cause households and tests near the edges to
be unallocated; overlaps will cause double-counting.

\textbf{2. Obtain source counts on their native geography.} For income,
retrieve ACS block-group estimates for the county of interest via the
Census API. The two variables needed are B19025\_001E (aggregate
household income) and B11001\_001E (number of households). Both must be
five-year estimates aligned with the same ACS vintage. For broadband,
download the relevant Ookla quarterly fixed-broadband tile from Ookla's
public S3 bucket. Clip the tile to your jurisdiction's bounding box to
reduce memory usage. Retain \texttt{tests} and \texttt{avg\_d\_kbps}.
Both datasets must be in the same CRS as your target geography before
redistribution.

\textbf{3. Redistribute the source counts to your target geography.}
Call \texttt{redistribute\_direct} (from the \texttt{sda\_areal} Python
package, or its R equivalent once available) with your source
GeoDataFrame and target GeoDataFrame. For ACS, redistribute
\texttt{agg\_income} and \texttt{households}. For Ookla, first construct
\texttt{d\_product\ =\ avg\_d\_kbps\ *\ tests}, then redistribute
\texttt{tests} and \texttt{d\_product}. Each call preserves the source
total within floating-point tolerance.

\textbf{4. Derive the intensive rates from the redistributed counts.}
Mean income is \texttt{agg\_income\_target\ /\ households\_target}.
Download speed in Mbps is
\texttt{(d\_product\_target\ /\ tests\_target)\ /\ 1000}. Neither
derived variable should be redistributed directly; the correct workflow
always redistributes the underlying counts and derives the rate
afterward. Attempting to redistribute mean income or average speed
directly will produce area-weighted results that are correct in only the
degenerate case where all source units are identical in density.

\textbf{5. Validate that totals are preserved.} Before proceeding to
analysis, compare the sum of redistributed households across all target
polygons to the sum of ACS households across all source block groups.
They should agree within a few percent. A larger discrepancy indicates a
geometry problem (gaps, overlaps, CRS mismatch) that will bias every
downstream estimate. Similarly check that redistributed test counts sum
to within a few percent of source tile test counts. Both checks are
implemented in the companion repository's \texttt{pipeline.validate}
module.

\textbf{6. Map, analyse, and interpret.} Join the income and speed
estimates to your target GeoDataFrame and export as GeoJSON, Parquet, or
whatever format your visualisation or analysis tool requires. The
companion repository builds all five result figures (choropleth maps,
scatter plot, and bivariate map) through a single command.

\section{Running the companion
pipeline}\label{running-the-companion-pipeline}

The full pipeline is executed with a single command from the repository
root:

\begin{Shaded}
\begin{Highlighting}[]
\ExtensionTok{uv}\NormalTok{ run python }\AttributeTok{{-}m}\NormalTok{ pipeline.run}
\end{Highlighting}
\end{Shaded}

This command fetches the Census and Ookla data, performs both
redistributions, validates totals, and writes the output GeoJSON. To
regenerate the result figures from the pipeline output:

\begin{Shaded}
\begin{Highlighting}[]
\ExtensionTok{uv}\NormalTok{ run python }\AttributeTok{{-}m}\NormalTok{ pipeline.build\_figures}
\end{Highlighting}
\end{Shaded}

Both commands require \texttt{uv} (the Python package manager), a Census
API key in the environment variable \texttt{CENSUS\_API\_KEY}, and an
internet connection. The repository README provides setup instructions.
All code is open-source and MIT-licensed.

\section{R implementation}\label{r-implementation}

The areal interpolation logic in the Python pipeline mirrors the
\texttt{sdc-redistribute} family of functions available in R. An R
implementation of the full Arlington pipeline is forthcoming, using
\texttt{tigris} for Census geometries, \texttt{censusapi} or
\texttt{tidycensus} for ACS data, \texttt{ooklaOpenDataR} for speed-test
tiles, and \texttt{areal} or a direct port of
\texttt{redistribute\_direct} for the interpolation step. Analysts
already working in R can follow the same six-step recipe using those
packages; the redistribution algebra is language-agnostic.

\section{Adapting to other variables}\label{adapting-to-other-variables}

The recipe generalises beyond income and broadband. Any ACS extensive
variable (counts of people, households, housing units, workers) can be
redistributed in step 3 by substituting the appropriate B-table
variables. Any Ookla variable that can be expressed as a count-weighted
product (upload speed follows the same \texttt{avg\_u\_kbps\ *\ tests}
construction as download speed) can be substituted in the same step.
Analysts working with other speed-test or coverage datasets will need to
identify the analogous extensive quantity (typically a count that, when
summed and divided into a product column, recovers the intensive rate of
interest) and follow the same pattern.

\bookmarksetup{startatroot}

\chapter{The Complete Pipeline Code}\label{sec-code}

This appendix reproduces the full, runnable Python pipeline behind the
worked example: the actual source, not a simplified sketch. It is
generated directly from the \texttt{pipeline/} package in the companion
repository, so it always matches the code that produced every figure and
number in this guide.

The pipeline is built on the
\href{https://pypi.org/project/sdc-redistribute/}{\texttt{sdc-redistribute}}
package and runs end to end with a single command:

\begin{Shaded}
\begin{Highlighting}[]
\ExtensionTok{uv}\NormalTok{ run python }\AttributeTok{{-}m}\NormalTok{ pipeline.run}
\end{Highlighting}
\end{Shaded}

The modules below appear in execution order. Figure generation
(\texttt{pipeline/style.py}, \texttt{pipeline/maps.py},
\texttt{pipeline/figures.py}) lives in the same repository and is
omitted here to keep the focus on the data method.

\section{Configuration}\label{configuration}

Central settings: the Arlington FIPS codes, the ACS count variables
(aggregate income and households), and the dataset paths every other
module reads.

\textbf{\texttt{pipeline/config.py}}

\begin{Shaded}
\begin{Highlighting}[]
\CommentTok{"""Central configuration: geographies, CRS, source identifiers, paths."""}
\ImportTok{from}\NormalTok{ pathlib }\ImportTok{import}\NormalTok{ Path}

\ImportTok{from}\NormalTok{ dotenv }\ImportTok{import}\NormalTok{ load\_dotenv}

\CommentTok{\# Load secrets (e.g. CENSUS\_API\_KEY) from a local .env if present. The file is}
\CommentTok{\# gitignored; loading it here makes the key available to any pipeline module.}
\NormalTok{load\_dotenv()}

\CommentTok{\# Arlington County, Virginia}
\NormalTok{STATE\_FIPS }\OperatorTok{=} \StringTok{"51"}
\NormalTok{COUNTY\_FIPS }\OperatorTok{=} \StringTok{"013"}
\NormalTok{ACS\_YEAR }\OperatorTok{=} \DecValTok{2021}          \CommentTok{\# 5{-}year ACS vintage used throughout}
\NormalTok{OOKLA\_YEAR }\OperatorTok{=} \DecValTok{2023}
\NormalTok{OOKLA\_QUARTER }\OperatorTok{=} \DecValTok{2}

\CommentTok{\# Area weighting is handled inside sdc{-}redistribute, which reprojects geographic}
\CommentTok{\# inputs to EPSG:3857 before computing intersection areas. At county scale the}
\CommentTok{\# choice of projected CRS is immaterial (\textless{}0.3\% area{-}ratio spread vs State Plane),}
\CommentTok{\# so the pipeline does not impose its own projected CRS.}

\CommentTok{\# ACS variables are EXTENSIVE COUNTS (never redistribute the median directly).}
\CommentTok{\# Mean household income is derived later as agg\_income / households.}
\NormalTok{ACS\_VARS }\OperatorTok{=}\NormalTok{ \{}
    \StringTok{"agg\_income"}\NormalTok{: }\StringTok{"B19025\_001"}\NormalTok{,   }\CommentTok{\# Aggregate household income (dollars)}
    \StringTok{"households"}\NormalTok{: }\StringTok{"B11001\_001"}\NormalTok{,   }\CommentTok{\# Total households (count)}
\NormalTok{\}}

\CommentTok{\# Ookla open data S3 parquet path template}
\NormalTok{OOKLA\_S3 }\OperatorTok{=}\NormalTok{ (}
    \StringTok{"s3://ookla{-}open{-}data/parquet/performance/type=fixed/"}
    \StringTok{"year=}\SpecialCharTok{\{year\}}\StringTok{/quarter=}\SpecialCharTok{\{quarter\}}\StringTok{/"}
    \StringTok{"}\SpecialCharTok{\{year\}}\StringTok{{-}}\SpecialCharTok{\{month:02d\}}\StringTok{{-}01\_performance\_fixed\_tiles.parquet"}
\NormalTok{)}

\CommentTok{\# Paths}
\NormalTok{ROOT }\OperatorTok{=}\NormalTok{ Path(}\VariableTok{\_\_file\_\_}\NormalTok{).resolve().parent.parent}
\NormalTok{DATA }\OperatorTok{=}\NormalTok{ ROOT }\OperatorTok{/} \StringTok{"data"}

\CommentTok{\# Source / output dataset paths}
\NormalTok{CIVIC\_ASSOC }\OperatorTok{=}\NormalTok{ DATA }\OperatorTok{/} \StringTok{"va013\_geo\_arl\_2021\_civic\_associations.geojson"}
\NormalTok{BLOCK\_GROUPS }\OperatorTok{=}\NormalTok{ DATA }\OperatorTok{/} \StringTok{"va013\_block\_groups.geojson"}          \CommentTok{\# produced by acquire\_geographies}
\NormalTok{BLOCKS }\OperatorTok{=}\NormalTok{ DATA }\OperatorTok{/} \StringTok{"va013\_geo\_blocks.geojson"}                  \CommentTok{\# existing}
\NormalTok{ACS\_COUNTS }\OperatorTok{=}\NormalTok{ DATA }\OperatorTok{/} \StringTok{"va013\_acs\_income\_counts.geojson"}       \CommentTok{\# produced by acquire\_acs}
\NormalTok{OOKLA\_TILES }\OperatorTok{=}\NormalTok{ DATA }\OperatorTok{/} \StringTok{"ookla\_tiles\_arlington.geojson"}        \CommentTok{\# existing/refreshable}
\NormalTok{CIVIC\_INCOME }\OperatorTok{=}\NormalTok{ DATA }\OperatorTok{/} \StringTok{"civic\_income.csv"}                    \CommentTok{\# produced by redistribute\_income}
\NormalTok{CIVIC\_BROADBAND }\OperatorTok{=}\NormalTok{ DATA }\OperatorTok{/} \StringTok{"civic\_broadband.csv"}             \CommentTok{\# produced by redistribute\_broadband}
\NormalTok{CIVIC\_COMBINED }\OperatorTok{=}\NormalTok{ DATA }\OperatorTok{/} \StringTok{"civic\_combined.geojson"}           \CommentTok{\# produced by combine}
\NormalTok{PARCELS }\OperatorTok{=}\NormalTok{ DATA }\OperatorTok{/} \StringTok{"va013\_parcels.geojson"}                   \CommentTok{\# produced by acquire\_parcels}
\NormalTok{CIVIC\_INCOME\_PARCELS }\OperatorTok{=}\NormalTok{ DATA }\OperatorTok{/} \StringTok{"civic\_income\_parcels.csv"}   \CommentTok{\# produced by redistribute\_income\_parcels}
\NormalTok{CIVIC\_INCOME\_COMPARISON }\OperatorTok{=}\NormalTok{ DATA }\OperatorTok{/} \StringTok{"civic\_income\_comparison.geojson"}  \CommentTok{\# produced by compare\_methods}
\end{Highlighting}
\end{Shaded}

\section{Acquire census geographies}\label{acquire-census-geographies}

Download Arlington's census block groups, the source geography for the
income redistribution.

\textbf{\texttt{pipeline/acquire\_geographies.py}}

\begin{Shaded}
\begin{Highlighting}[]
\CommentTok{"""Acquire Arlington census block{-}group geometries via pygris."""}
\ImportTok{from}\NormalTok{ \_\_future\_\_ }\ImportTok{import}\NormalTok{ annotations}

\ImportTok{import}\NormalTok{ geopandas }\ImportTok{as}\NormalTok{ gpd}
\ImportTok{from}\NormalTok{ pygris }\ImportTok{import}\NormalTok{ block\_groups}

\ImportTok{from}\NormalTok{ pipeline }\ImportTok{import}\NormalTok{ config}


\KeywordTok{def}\NormalTok{ fetch\_block\_groups() }\OperatorTok{{-}\textgreater{}}\NormalTok{ gpd.GeoDataFrame:}
    \CommentTok{"""Download Arlington County (VA/013) block groups for the ACS vintage."""}
\NormalTok{    bg }\OperatorTok{=}\NormalTok{ block\_groups(}
\NormalTok{        state}\OperatorTok{=}\NormalTok{config.STATE\_FIPS,}
\NormalTok{        county}\OperatorTok{=}\NormalTok{config.COUNTY\_FIPS,}
\NormalTok{        year}\OperatorTok{=}\NormalTok{config.ACS\_YEAR,}
\NormalTok{        cb}\OperatorTok{=}\VariableTok{True}\NormalTok{,}
\NormalTok{    )}
    \ControlFlowTok{return}\NormalTok{ bg}


\KeywordTok{def}\NormalTok{ main() }\OperatorTok{{-}\textgreater{}} \VariableTok{None}\NormalTok{:}
\NormalTok{    bg }\OperatorTok{=}\NormalTok{ fetch\_block\_groups()}
\NormalTok{    bg.to\_file(config.BLOCK\_GROUPS, driver}\OperatorTok{=}\StringTok{"GeoJSON"}\NormalTok{)}
    \BuiltInTok{print}\NormalTok{(}\SpecialStringTok{f"Wrote }\SpecialCharTok{\{}\BuiltInTok{len}\NormalTok{(bg)}\SpecialCharTok{\}}\SpecialStringTok{ block groups → }\SpecialCharTok{\{}\NormalTok{config}\SpecialCharTok{.}\NormalTok{BLOCK\_GROUPS}\SpecialCharTok{\}}\SpecialStringTok{"}\NormalTok{)}


\ControlFlowTok{if} \VariableTok{\_\_name\_\_} \OperatorTok{==} \StringTok{"\_\_main\_\_"}\NormalTok{:}
\NormalTok{    main()}
\end{Highlighting}
\end{Shaded}

\section{Acquire ACS income counts}\label{acquire-acs-income-counts}

Pull aggregate household income and household counts from the Census
API. These are \emph{counts} (extensive measures); mean income is
derived later.

\textbf{\texttt{pipeline/acquire\_acs.py}}

\begin{Shaded}
\begin{Highlighting}[]
\CommentTok{"""Acquire Arlington ACS aggregate income + household counts (block group)."""}
\ImportTok{from}\NormalTok{ \_\_future\_\_ }\ImportTok{import}\NormalTok{ annotations}

\ImportTok{import}\NormalTok{ os}

\ImportTok{import}\NormalTok{ geopandas }\ImportTok{as}\NormalTok{ gpd}
\ImportTok{import}\NormalTok{ pandas }\ImportTok{as}\NormalTok{ pd}
\ImportTok{from}\NormalTok{ census }\ImportTok{import}\NormalTok{ Census}

\ImportTok{from}\NormalTok{ pipeline }\ImportTok{import}\NormalTok{ config}
\ImportTok{from}\NormalTok{ pipeline.acquire\_geographies }\ImportTok{import}\NormalTok{ fetch\_block\_groups}


\KeywordTok{def}\NormalTok{ \_sanitize\_counts(df: pd.DataFrame) }\OperatorTok{{-}\textgreater{}}\NormalTok{ pd.DataFrame:}
    \CommentTok{"""Replace Census negative sentinels (e.g. {-}666666666) with 0.}

\CommentTok{    Aggregate income and household counts are non{-}negative by definition; the}
\CommentTok{    Census API returns large negative sentinel codes for suppressed or}
\CommentTok{    not{-}applicable cells (typically zero{-}household block groups). Left in place,}
\CommentTok{    these sentinels would propagate through area{-}weighted redistribution and}
\CommentTok{    produce negative incomes. Treating them as 0 is correct for the}
\CommentTok{    zero{-}household areas they mark.}
\CommentTok{    """}
\NormalTok{    out }\OperatorTok{=}\NormalTok{ df.copy()}
    \ControlFlowTok{for}\NormalTok{ col }\KeywordTok{in}\NormalTok{ (}\StringTok{"agg\_income"}\NormalTok{, }\StringTok{"households"}\NormalTok{):}
\NormalTok{        out[col] }\OperatorTok{=}\NormalTok{ out[col].mask(out[col] }\OperatorTok{\textless{}} \DecValTok{0}\NormalTok{, }\DecValTok{0}\NormalTok{)}
    \ControlFlowTok{return}\NormalTok{ out}


\KeywordTok{def}\NormalTok{ fetch\_acs\_counts(api\_key: }\BuiltInTok{str} \OperatorTok{|} \VariableTok{None} \OperatorTok{=} \VariableTok{None}\NormalTok{) }\OperatorTok{{-}\textgreater{}}\NormalTok{ pd.DataFrame:}
    \CommentTok{"""Return a DataFrame: GEOID, agg\_income, households for Arlington block groups."""}
\NormalTok{    key }\OperatorTok{=}\NormalTok{ api\_key }\KeywordTok{or}\NormalTok{ os.environ.get(}\StringTok{"CENSUS\_API\_KEY"}\NormalTok{)}
\NormalTok{    c }\OperatorTok{=}\NormalTok{ Census(key, year}\OperatorTok{=}\NormalTok{config.ACS\_YEAR)}
    \CommentTok{\# The census library needs the API estimate suffix "E" on each variable}
    \CommentTok{\# (e.g. B19025\_001E), unlike tidycensus which strips it.}
\NormalTok{    agg\_var }\OperatorTok{=}\NormalTok{ config.ACS\_VARS[}\StringTok{"agg\_income"}\NormalTok{] }\OperatorTok{+} \StringTok{"E"}
\NormalTok{    hh\_var }\OperatorTok{=}\NormalTok{ config.ACS\_VARS[}\StringTok{"households"}\NormalTok{] }\OperatorTok{+} \StringTok{"E"}
\NormalTok{    rows }\OperatorTok{=}\NormalTok{ c.acs5.state\_county\_blockgroup(}
\NormalTok{        fields}\OperatorTok{=}\NormalTok{(}\StringTok{"NAME"}\NormalTok{, agg\_var, hh\_var),}
\NormalTok{        state\_fips}\OperatorTok{=}\NormalTok{config.STATE\_FIPS,}
\NormalTok{        county\_fips}\OperatorTok{=}\NormalTok{config.COUNTY\_FIPS,}
\NormalTok{        blockgroup}\OperatorTok{=}\NormalTok{Census.ALL,}
\NormalTok{        tract}\OperatorTok{=}\NormalTok{Census.ALL,}
\NormalTok{    )}
\NormalTok{    df }\OperatorTok{=}\NormalTok{ pd.DataFrame(rows)}
\NormalTok{    df[}\StringTok{"GEOID"}\NormalTok{] }\OperatorTok{=}\NormalTok{ df[}\StringTok{"state"}\NormalTok{] }\OperatorTok{+}\NormalTok{ df[}\StringTok{"county"}\NormalTok{] }\OperatorTok{+}\NormalTok{ df[}\StringTok{"tract"}\NormalTok{] }\OperatorTok{+}\NormalTok{ df[}\StringTok{"block group"}\NormalTok{]}
\NormalTok{    df }\OperatorTok{=}\NormalTok{ df.rename(columns}\OperatorTok{=}\NormalTok{\{agg\_var: }\StringTok{"agg\_income"}\NormalTok{, hh\_var: }\StringTok{"households"}\NormalTok{\})}
\NormalTok{    df }\OperatorTok{=}\NormalTok{ \_sanitize\_counts(df)}
    \ControlFlowTok{return}\NormalTok{ df[[}\StringTok{"GEOID"}\NormalTok{, }\StringTok{"agg\_income"}\NormalTok{, }\StringTok{"households"}\NormalTok{]]}


\KeywordTok{def}\NormalTok{ main() }\OperatorTok{{-}\textgreater{}} \VariableTok{None}\NormalTok{:}
\NormalTok{    counts }\OperatorTok{=}\NormalTok{ fetch\_acs\_counts()}
\NormalTok{    bg }\OperatorTok{=}\NormalTok{ fetch\_block\_groups()[[}\StringTok{"GEOID"}\NormalTok{, }\StringTok{"geometry"}\NormalTok{]]}
\NormalTok{    merged }\OperatorTok{=}\NormalTok{ bg.merge(counts, on}\OperatorTok{=}\StringTok{"GEOID"}\NormalTok{, how}\OperatorTok{=}\StringTok{"left"}\NormalTok{)}
\NormalTok{    gdf }\OperatorTok{=}\NormalTok{ gpd.GeoDataFrame(merged, geometry}\OperatorTok{=}\StringTok{"geometry"}\NormalTok{, crs}\OperatorTok{=}\NormalTok{bg.crs)}
\NormalTok{    gdf.to\_file(config.ACS\_COUNTS, driver}\OperatorTok{=}\StringTok{"GeoJSON"}\NormalTok{)}
    \BuiltInTok{print}\NormalTok{(}\SpecialStringTok{f"Wrote }\SpecialCharTok{\{}\BuiltInTok{len}\NormalTok{(gdf)}\SpecialCharTok{\}}\SpecialStringTok{ block groups with income counts → }\SpecialCharTok{\{}\NormalTok{config}\SpecialCharTok{.}\NormalTok{ACS\_COUNTS}\SpecialCharTok{\}}\SpecialStringTok{"}\NormalTok{)}


\ControlFlowTok{if} \VariableTok{\_\_name\_\_} \OperatorTok{==} \StringTok{"\_\_main\_\_"}\NormalTok{:}
\NormalTok{    main()}
\end{Highlighting}
\end{Shaded}

\section{Acquire Ookla broadband
tiles}\label{acquire-ookla-broadband-tiles}

Read the Ookla open-data broadband tiles from S3 and clip them to
Arlington.

\textbf{\texttt{pipeline/acquire\_ookla.py}}

\begin{Shaded}
\begin{Highlighting}[]
\CommentTok{"""Acquire Ookla fixed{-}broadband tiles for Arlington from S3 open data."""}
\ImportTok{from}\NormalTok{ \_\_future\_\_ }\ImportTok{import}\NormalTok{ annotations}

\ImportTok{import}\NormalTok{ geopandas }\ImportTok{as}\NormalTok{ gpd}
\ImportTok{import}\NormalTok{ pandas }\ImportTok{as}\NormalTok{ pd}
\ImportTok{from}\NormalTok{ shapely }\ImportTok{import}\NormalTok{ wkt}

\ImportTok{from}\NormalTok{ pipeline }\ImportTok{import}\NormalTok{ config}
\ImportTok{from}\NormalTok{ pipeline.acquire\_geographies }\ImportTok{import}\NormalTok{ fetch\_block\_groups}


\KeywordTok{def}\NormalTok{ \_s3\_path() }\OperatorTok{{-}\textgreater{}} \BuiltInTok{str}\NormalTok{:}
\NormalTok{    month }\OperatorTok{=}\NormalTok{ config.OOKLA\_QUARTER }\OperatorTok{*} \DecValTok{3} \OperatorTok{{-}} \DecValTok{2}  \CommentTok{\# Q1{-}\textgreater{}01, Q2{-}\textgreater{}04, Q3{-}\textgreater{}07, Q4{-}\textgreater{}10}
    \ControlFlowTok{return}\NormalTok{ config.OOKLA\_S3.}\BuiltInTok{format}\NormalTok{(}
\NormalTok{        year}\OperatorTok{=}\NormalTok{config.OOKLA\_YEAR, quarter}\OperatorTok{=}\NormalTok{config.OOKLA\_QUARTER, month}\OperatorTok{=}\NormalTok{month}
\NormalTok{    )}


\KeywordTok{def}\NormalTok{ fetch\_ookla\_tiles() }\OperatorTok{{-}\textgreater{}}\NormalTok{ gpd.GeoDataFrame:}
    \CommentTok{"""Download and bbox{-}filter Ookla tiles to Arlington County."""}
\NormalTok{    df }\OperatorTok{=}\NormalTok{ pd.read\_parquet(\_s3\_path(), storage\_options}\OperatorTok{=}\NormalTok{\{}\StringTok{"anon"}\NormalTok{: }\VariableTok{True}\NormalTok{\})}
\NormalTok{    df[}\StringTok{"geometry"}\NormalTok{] }\OperatorTok{=}\NormalTok{ df[}\StringTok{"tile"}\NormalTok{].}\BuiltInTok{apply}\NormalTok{(wkt.loads)}
\NormalTok{    tiles }\OperatorTok{=}\NormalTok{ gpd.GeoDataFrame(df, geometry}\OperatorTok{=}\StringTok{"geometry"}\NormalTok{, crs}\OperatorTok{=}\StringTok{"EPSG:4326"}\NormalTok{)}

\NormalTok{    bbox }\OperatorTok{=}\NormalTok{ fetch\_block\_groups().to\_crs(}\StringTok{"EPSG:4326"}\NormalTok{).total\_bounds  }\CommentTok{\# minx,miny,maxx,maxy}
\NormalTok{    minx, miny, maxx, maxy }\OperatorTok{=}\NormalTok{ bbox}
    \ControlFlowTok{return}\NormalTok{ tiles.cx[minx:maxx, miny:maxy].copy()}


\KeywordTok{def}\NormalTok{ main() }\OperatorTok{{-}\textgreater{}} \VariableTok{None}\NormalTok{:}
\NormalTok{    tiles }\OperatorTok{=}\NormalTok{ fetch\_ookla\_tiles()}
\NormalTok{    cols }\OperatorTok{=}\NormalTok{ [}\StringTok{"quadkey"}\NormalTok{, }\StringTok{"avg\_d\_kbps"}\NormalTok{, }\StringTok{"avg\_u\_kbps"}\NormalTok{, }\StringTok{"avg\_lat\_ms"}\NormalTok{, }\StringTok{"tests"}\NormalTok{, }\StringTok{"devices"}\NormalTok{, }\StringTok{"geometry"}\NormalTok{]}
\NormalTok{    tiles[cols].to\_file(config.OOKLA\_TILES, driver}\OperatorTok{=}\StringTok{"GeoJSON"}\NormalTok{)}
    \BuiltInTok{print}\NormalTok{(}\SpecialStringTok{f"Wrote }\SpecialCharTok{\{}\BuiltInTok{len}\NormalTok{(tiles)}\SpecialCharTok{\}}\SpecialStringTok{ Ookla tiles → }\SpecialCharTok{\{}\NormalTok{config}\SpecialCharTok{.}\NormalTok{OOKLA\_TILES}\SpecialCharTok{\}}\SpecialStringTok{"}\NormalTok{)}


\ControlFlowTok{if} \VariableTok{\_\_name\_\_} \OperatorTok{==} \StringTok{"\_\_main\_\_"}\NormalTok{:}
\NormalTok{    main()}
\end{Highlighting}
\end{Shaded}

\section{Redistribute income}\label{redistribute-income}

Area-weight the income and household counts onto civic associations with
\texttt{sdc-redistribute}, then derive mean household income.

\textbf{\texttt{pipeline/redistribute\_income.py}}

\begin{Shaded}
\begin{Highlighting}[]
\CommentTok{"""Redistribute ACS income counts to civic associations and derive mean income."""}
\ImportTok{from}\NormalTok{ \_\_future\_\_ }\ImportTok{import}\NormalTok{ annotations}

\ImportTok{from}\NormalTok{ pathlib }\ImportTok{import}\NormalTok{ Path}

\ImportTok{import}\NormalTok{ geopandas }\ImportTok{as}\NormalTok{ gpd}
\ImportTok{import}\NormalTok{ pandas }\ImportTok{as}\NormalTok{ pd}
\ImportTok{from}\NormalTok{ sdc\_redistribute }\ImportTok{import}\NormalTok{ redistribute\_direct}

\ImportTok{from}\NormalTok{ pipeline }\ImportTok{import}\NormalTok{ config}


\KeywordTok{def}\NormalTok{ redistribute\_income(}
\NormalTok{    counts\_long: pd.DataFrame, source\_geo: Path, target\_geo: Path}
\NormalTok{) }\OperatorTok{{-}\textgreater{}}\NormalTok{ pd.DataFrame:}
    \CommentTok{"""Return per{-}civic{-}association agg\_income, households, and derived mean\_income.}

\CommentTok{    \textasciigrave{}\textasciigrave{}counts\_long\textasciigrave{}\textasciigrave{} is long format (geoid, year, measure, value) with measures}
\CommentTok{    \textasciigrave{}\textasciigrave{}agg\_income\textasciigrave{}\textasciigrave{} and \textasciigrave{}\textasciigrave{}households\textasciigrave{}\textasciigrave{}.}
\CommentTok{    """}
\NormalTok{    result }\OperatorTok{=}\NormalTok{ redistribute\_direct(}
\NormalTok{        source\_df}\OperatorTok{=}\NormalTok{counts\_long,}
\NormalTok{        source\_geo}\OperatorTok{=}\NormalTok{source\_geo,}
\NormalTok{        target\_geos}\OperatorTok{=}\NormalTok{\{}\StringTok{"civic\_association"}\NormalTok{: target\_geo\},}
\NormalTok{        count\_cols}\OperatorTok{=}\NormalTok{[}\StringTok{"agg\_income"}\NormalTok{, }\StringTok{"households"}\NormalTok{],}
\NormalTok{        source\_id}\OperatorTok{=}\StringTok{"geoid"}\NormalTok{,}
\NormalTok{    )}
    \CommentTok{\# result is long with measures agg\_income\_direct / households\_direct}
\NormalTok{    wide }\OperatorTok{=}\NormalTok{ result.pivot\_table(}
\NormalTok{        index}\OperatorTok{=}\StringTok{"geoid"}\NormalTok{, columns}\OperatorTok{=}\StringTok{"measure"}\NormalTok{, values}\OperatorTok{=}\StringTok{"value"}\NormalTok{, aggfunc}\OperatorTok{=}\StringTok{"first"}
\NormalTok{    ).reset\_index()}
\NormalTok{    wide.columns.name }\OperatorTok{=} \VariableTok{None}
\NormalTok{    wide }\OperatorTok{=}\NormalTok{ wide.rename(}
\NormalTok{        columns}\OperatorTok{=}\NormalTok{\{}\StringTok{"agg\_income\_direct"}\NormalTok{: }\StringTok{"agg\_income"}\NormalTok{, }\StringTok{"households\_direct"}\NormalTok{: }\StringTok{"households"}\NormalTok{\}}
\NormalTok{    )}
    \CommentTok{\# Mean household income is an INTENSIVE quantity derived from two counts.}
\NormalTok{    wide[}\StringTok{"mean\_income"}\NormalTok{] }\OperatorTok{=}\NormalTok{ wide[}\StringTok{"agg\_income"}\NormalTok{] }\OperatorTok{/}\NormalTok{ wide[}\StringTok{"households"}\NormalTok{].where(}
\NormalTok{        wide[}\StringTok{"households"}\NormalTok{] }\OperatorTok{\textgreater{}} \DecValTok{0}
\NormalTok{    )}
    \ControlFlowTok{return}\NormalTok{ wide[[}\StringTok{"geoid"}\NormalTok{, }\StringTok{"agg\_income"}\NormalTok{, }\StringTok{"households"}\NormalTok{, }\StringTok{"mean\_income"}\NormalTok{]]}


\KeywordTok{def}\NormalTok{ main() }\OperatorTok{{-}\textgreater{}} \VariableTok{None}\NormalTok{:}
    \CommentTok{\# ACS counts GeoJSON uses uppercase GEOID; rename to the lowercase \textasciigrave{}geoid\textasciigrave{}}
    \CommentTok{\# that redistribute\_direct expects, for both the long{-}format counts and the}
    \CommentTok{\# source geometry. Write a geoid{-}keyed source file for redistribute\_direct.}
\NormalTok{    bg }\OperatorTok{=}\NormalTok{ gpd.read\_file(config.ACS\_COUNTS).rename(columns}\OperatorTok{=}\NormalTok{\{}\StringTok{"GEOID"}\NormalTok{: }\StringTok{"geoid"}\NormalTok{\})}
    \BuiltInTok{long} \OperatorTok{=}\NormalTok{ bg[[}\StringTok{"geoid"}\NormalTok{, }\StringTok{"agg\_income"}\NormalTok{, }\StringTok{"households"}\NormalTok{]].melt(}
\NormalTok{        id\_vars}\OperatorTok{=}\StringTok{"geoid"}\NormalTok{, var\_name}\OperatorTok{=}\StringTok{"measure"}\NormalTok{, value\_name}\OperatorTok{=}\StringTok{"value"}
\NormalTok{    )}
    \BuiltInTok{long}\NormalTok{[}\StringTok{"year"}\NormalTok{] }\OperatorTok{=}\NormalTok{ config.ACS\_YEAR}

\NormalTok{    src\_path }\OperatorTok{=}\NormalTok{ config.DATA }\OperatorTok{/} \StringTok{"\_acs\_src.geojson"}
\NormalTok{    bg[[}\StringTok{"geoid"}\NormalTok{, }\StringTok{"geometry"}\NormalTok{]].to\_file(src\_path, driver}\OperatorTok{=}\StringTok{"GeoJSON"}\NormalTok{)}
\NormalTok{    out }\OperatorTok{=}\NormalTok{ redistribute\_income(}\BuiltInTok{long}\NormalTok{, src\_path, config.CIVIC\_ASSOC)}
\NormalTok{    src\_path.unlink(missing\_ok}\OperatorTok{=}\VariableTok{True}\NormalTok{)}

\NormalTok{    out.to\_csv(config.CIVIC\_INCOME, index}\OperatorTok{=}\VariableTok{False}\NormalTok{)}
    \BuiltInTok{print}\NormalTok{(}\SpecialStringTok{f"Wrote }\SpecialCharTok{\{}\BuiltInTok{len}\NormalTok{(out)}\SpecialCharTok{\}}\SpecialStringTok{ civic{-}association income rows → }\SpecialCharTok{\{}\NormalTok{config}\SpecialCharTok{.}\NormalTok{CIVIC\_INCOME}\SpecialCharTok{\}}\SpecialStringTok{"}\NormalTok{)}


\ControlFlowTok{if} \VariableTok{\_\_name\_\_} \OperatorTok{==} \StringTok{"\_\_main\_\_"}\NormalTok{:}
\NormalTok{    main()}
\end{Highlighting}
\end{Shaded}

\section{Redistribute broadband}\label{redistribute-broadband}

Redistribute test counts and the speed×tests product, then derive the
count-weighted mean download and upload speeds.

\textbf{\texttt{pipeline/redistribute\_broadband.py}}

\begin{Shaded}
\begin{Highlighting}[]
\CommentTok{"""Redistribute Ookla broadband to civic associations; derive mean speeds."""}
\ImportTok{from}\NormalTok{ \_\_future\_\_ }\ImportTok{import}\NormalTok{ annotations}

\ImportTok{from}\NormalTok{ pathlib }\ImportTok{import}\NormalTok{ Path}

\ImportTok{import}\NormalTok{ geopandas }\ImportTok{as}\NormalTok{ gpd}
\ImportTok{import}\NormalTok{ pandas }\ImportTok{as}\NormalTok{ pd}
\ImportTok{from}\NormalTok{ sdc\_redistribute }\ImportTok{import}\NormalTok{ redistribute\_direct}

\ImportTok{from}\NormalTok{ pipeline }\ImportTok{import}\NormalTok{ config}

\NormalTok{MIN\_TESTS }\OperatorTok{=} \DecValTok{5}  \CommentTok{\# statistical{-}reliability filter on source tiles}


\KeywordTok{def}\NormalTok{ redistribute\_broadband(}
\NormalTok{    counts\_long: pd.DataFrame, source\_geo: Path, target\_geo: Path}
\NormalTok{) }\OperatorTok{{-}\textgreater{}}\NormalTok{ pd.DataFrame:}
    \CommentTok{"""Return per{-}civic{-}association tests and derived download/upload Mbps."""}
\NormalTok{    result }\OperatorTok{=}\NormalTok{ redistribute\_direct(}
\NormalTok{        source\_df}\OperatorTok{=}\NormalTok{counts\_long,}
\NormalTok{        source\_geo}\OperatorTok{=}\NormalTok{source\_geo,}
\NormalTok{        target\_geos}\OperatorTok{=}\NormalTok{\{}\StringTok{"civic\_association"}\NormalTok{: target\_geo\},}
\NormalTok{        count\_cols}\OperatorTok{=}\NormalTok{[}\StringTok{"tests"}\NormalTok{, }\StringTok{"d\_product"}\NormalTok{, }\StringTok{"u\_product"}\NormalTok{],}
\NormalTok{        source\_id}\OperatorTok{=}\StringTok{"geoid"}\NormalTok{,}
\NormalTok{    )}
\NormalTok{    wide }\OperatorTok{=}\NormalTok{ result.pivot\_table(}
\NormalTok{        index}\OperatorTok{=}\StringTok{"geoid"}\NormalTok{, columns}\OperatorTok{=}\StringTok{"measure"}\NormalTok{, values}\OperatorTok{=}\StringTok{"value"}\NormalTok{, aggfunc}\OperatorTok{=}\StringTok{"first"}
\NormalTok{    ).reset\_index()}
\NormalTok{    wide.columns.name }\OperatorTok{=} \VariableTok{None}
\NormalTok{    wide }\OperatorTok{=}\NormalTok{ wide.rename(}
\NormalTok{        columns}\OperatorTok{=}\NormalTok{\{}
            \StringTok{"tests\_direct"}\NormalTok{: }\StringTok{"tests"}\NormalTok{,}
            \StringTok{"d\_product\_direct"}\NormalTok{: }\StringTok{"d\_product"}\NormalTok{,}
            \StringTok{"u\_product\_direct"}\NormalTok{: }\StringTok{"u\_product"}\NormalTok{,}
\NormalTok{        \}}
\NormalTok{    )}
\NormalTok{    valid }\OperatorTok{=}\NormalTok{ wide[}\StringTok{"tests"}\NormalTok{] }\OperatorTok{\textgreater{}} \DecValTok{0}
\NormalTok{    wide[}\StringTok{"download\_mbps"}\NormalTok{] }\OperatorTok{=}\NormalTok{ (wide[}\StringTok{"d\_product"}\NormalTok{] }\OperatorTok{/}\NormalTok{ wide[}\StringTok{"tests"}\NormalTok{].where(valid)) }\OperatorTok{/} \DecValTok{1000}
\NormalTok{    wide[}\StringTok{"upload\_mbps"}\NormalTok{] }\OperatorTok{=}\NormalTok{ (wide[}\StringTok{"u\_product"}\NormalTok{] }\OperatorTok{/}\NormalTok{ wide[}\StringTok{"tests"}\NormalTok{].where(valid)) }\OperatorTok{/} \DecValTok{1000}
    \ControlFlowTok{return}\NormalTok{ wide[[}\StringTok{"geoid"}\NormalTok{, }\StringTok{"tests"}\NormalTok{, }\StringTok{"download\_mbps"}\NormalTok{, }\StringTok{"upload\_mbps"}\NormalTok{]]}


\KeywordTok{def}\NormalTok{ build\_counts\_long(tiles: gpd.GeoDataFrame) }\OperatorTok{{-}\textgreater{}}\NormalTok{ pd.DataFrame:}
    \CommentTok{"""Build the long{-}format count frame from raw Ookla tiles."""}
\NormalTok{    t }\OperatorTok{=}\NormalTok{ tiles[tiles[}\StringTok{"tests"}\NormalTok{] }\OperatorTok{\textgreater{}=}\NormalTok{ MIN\_TESTS].copy()}
\NormalTok{    t[}\StringTok{"d\_product"}\NormalTok{] }\OperatorTok{=}\NormalTok{ t[}\StringTok{"avg\_d\_kbps"}\NormalTok{] }\OperatorTok{*}\NormalTok{ t[}\StringTok{"tests"}\NormalTok{]}
\NormalTok{    t[}\StringTok{"u\_product"}\NormalTok{] }\OperatorTok{=}\NormalTok{ t[}\StringTok{"avg\_u\_kbps"}\NormalTok{] }\OperatorTok{*}\NormalTok{ t[}\StringTok{"tests"}\NormalTok{]}
    \BuiltInTok{long} \OperatorTok{=}\NormalTok{ t[[}\StringTok{"quadkey"}\NormalTok{, }\StringTok{"tests"}\NormalTok{, }\StringTok{"d\_product"}\NormalTok{, }\StringTok{"u\_product"}\NormalTok{]].melt(}
\NormalTok{        id\_vars}\OperatorTok{=}\StringTok{"quadkey"}\NormalTok{, var\_name}\OperatorTok{=}\StringTok{"measure"}\NormalTok{, value\_name}\OperatorTok{=}\StringTok{"value"}
\NormalTok{    )}
    \BuiltInTok{long} \OperatorTok{=} \BuiltInTok{long}\NormalTok{.rename(columns}\OperatorTok{=}\NormalTok{\{}\StringTok{"quadkey"}\NormalTok{: }\StringTok{"geoid"}\NormalTok{\})}
    \BuiltInTok{long}\NormalTok{[}\StringTok{"year"}\NormalTok{] }\OperatorTok{=}\NormalTok{ config.OOKLA\_YEAR}
    \ControlFlowTok{return} \BuiltInTok{long}


\KeywordTok{def}\NormalTok{ main() }\OperatorTok{{-}\textgreater{}} \VariableTok{None}\NormalTok{:}
\NormalTok{    tiles }\OperatorTok{=}\NormalTok{ gpd.read\_file(config.OOKLA\_TILES)}
\NormalTok{    tiles[}\StringTok{"geoid"}\NormalTok{] }\OperatorTok{=}\NormalTok{ tiles[}\StringTok{"quadkey"}\NormalTok{].astype(}\BuiltInTok{str}\NormalTok{)}
    \BuiltInTok{long} \OperatorTok{=}\NormalTok{ build\_counts\_long(tiles)}
    \CommentTok{\# Write a geoid{-}keyed source file for redistribute\_direct}
\NormalTok{    src }\OperatorTok{=}\NormalTok{ tiles[[}\StringTok{"geoid"}\NormalTok{, }\StringTok{"geometry"}\NormalTok{]].copy()}
\NormalTok{    src\_path }\OperatorTok{=}\NormalTok{ config.DATA }\OperatorTok{/} \StringTok{"\_ookla\_src.geojson"}
\NormalTok{    src.to\_file(src\_path, driver}\OperatorTok{=}\StringTok{"GeoJSON"}\NormalTok{)}
\NormalTok{    out }\OperatorTok{=}\NormalTok{ redistribute\_broadband(}\BuiltInTok{long}\NormalTok{, src\_path, config.CIVIC\_ASSOC)}
\NormalTok{    out.to\_csv(config.CIVIC\_BROADBAND, index}\OperatorTok{=}\VariableTok{False}\NormalTok{)}
\NormalTok{    src\_path.unlink(missing\_ok}\OperatorTok{=}\VariableTok{True}\NormalTok{)}
    \BuiltInTok{print}\NormalTok{(}\SpecialStringTok{f"Wrote }\SpecialCharTok{\{}\BuiltInTok{len}\NormalTok{(out)}\SpecialCharTok{\}}\SpecialStringTok{ civic{-}association broadband rows → }\SpecialCharTok{\{}\NormalTok{config}\SpecialCharTok{.}\NormalTok{CIVIC\_BROADBAND}\SpecialCharTok{\}}\SpecialStringTok{"}\NormalTok{)}


\ControlFlowTok{if} \VariableTok{\_\_name\_\_} \OperatorTok{==} \StringTok{"\_\_main\_\_"}\NormalTok{:}
\NormalTok{    main()}
\end{Highlighting}
\end{Shaded}

\section{Combine and derive metrics}\label{combine-and-derive-metrics}

Join the redistributed income and broadband measures onto the civic
association polygons and compute the income-to-speed ratio and the
bivariate income×speed classes.

\textbf{\texttt{pipeline/combine.py}}

\begin{Shaded}
\begin{Highlighting}[]
\CommentTok{"""Combine income + broadband and compute derived policy metrics."""}
\ImportTok{from}\NormalTok{ \_\_future\_\_ }\ImportTok{import}\NormalTok{ annotations}

\ImportTok{import}\NormalTok{ geopandas }\ImportTok{as}\NormalTok{ gpd}
\ImportTok{import}\NormalTok{ pandas }\ImportTok{as}\NormalTok{ pd}

\ImportTok{from}\NormalTok{ pipeline }\ImportTok{import}\NormalTok{ config}


\KeywordTok{def}\NormalTok{ \_tercile(series: pd.Series) }\OperatorTok{{-}\textgreater{}}\NormalTok{ pd.Series:}
    \CommentTok{"""Return 1/2/3 tercile labels by rank (robust to ties and NaN)."""}
\NormalTok{    ranks }\OperatorTok{=}\NormalTok{ series.rank(method}\OperatorTok{=}\StringTok{"first"}\NormalTok{)}
    \ControlFlowTok{return}\NormalTok{ pd.qcut(ranks, }\DecValTok{3}\NormalTok{, labels}\OperatorTok{=}\NormalTok{[}\DecValTok{1}\NormalTok{, }\DecValTok{2}\NormalTok{, }\DecValTok{3}\NormalTok{]).astype(}\StringTok{"Int64"}\NormalTok{)}


\KeywordTok{def}\NormalTok{ add\_derived\_metrics(df: pd.DataFrame) }\OperatorTok{{-}\textgreater{}}\NormalTok{ pd.DataFrame:}
    \CommentTok{"""Add income\_speed\_ratio and bivariate\_class (income{-}speed terciles)."""}
\NormalTok{    out }\OperatorTok{=}\NormalTok{ df.copy()}
\NormalTok{    out[}\StringTok{"income\_speed\_ratio"}\NormalTok{] }\OperatorTok{=}\NormalTok{ out[}\StringTok{"mean\_income"}\NormalTok{] }\OperatorTok{/}\NormalTok{ out[}\StringTok{"download\_mbps"}\NormalTok{].where(}
\NormalTok{        out[}\StringTok{"download\_mbps"}\NormalTok{] }\OperatorTok{\textgreater{}} \DecValTok{0}
\NormalTok{    )}
\NormalTok{    inc }\OperatorTok{=}\NormalTok{ \_tercile(out[}\StringTok{"mean\_income"}\NormalTok{])}
\NormalTok{    spd }\OperatorTok{=}\NormalTok{ \_tercile(out[}\StringTok{"download\_mbps"}\NormalTok{])}
\NormalTok{    out[}\StringTok{"bivariate\_class"}\NormalTok{] }\OperatorTok{=}\NormalTok{ inc.astype(}\BuiltInTok{str}\NormalTok{) }\OperatorTok{+} \StringTok{"{-}"} \OperatorTok{+}\NormalTok{ spd.astype(}\BuiltInTok{str}\NormalTok{)}
    \ControlFlowTok{return}\NormalTok{ out}


\KeywordTok{def}\NormalTok{ main() }\OperatorTok{{-}\textgreater{}} \VariableTok{None}\NormalTok{:}
\NormalTok{    civ }\OperatorTok{=}\NormalTok{ gpd.read\_file(config.CIVIC\_ASSOC)[[}\StringTok{"geoid"}\NormalTok{, }\StringTok{"region\_name"}\NormalTok{, }\StringTok{"geometry"}\NormalTok{]]}
\NormalTok{    income }\OperatorTok{=}\NormalTok{ pd.read\_csv(config.CIVIC\_INCOME, dtype}\OperatorTok{=}\NormalTok{\{}\StringTok{"geoid"}\NormalTok{: }\BuiltInTok{str}\NormalTok{\})}
\NormalTok{    broadband }\OperatorTok{=}\NormalTok{ pd.read\_csv(config.CIVIC\_BROADBAND, dtype}\OperatorTok{=}\NormalTok{\{}\StringTok{"geoid"}\NormalTok{: }\BuiltInTok{str}\NormalTok{\})}
\NormalTok{    merged }\OperatorTok{=}\NormalTok{ civ.merge(income, on}\OperatorTok{=}\StringTok{"geoid"}\NormalTok{, how}\OperatorTok{=}\StringTok{"left"}\NormalTok{).merge(}
\NormalTok{        broadband, on}\OperatorTok{=}\StringTok{"geoid"}\NormalTok{, how}\OperatorTok{=}\StringTok{"left"}
\NormalTok{    )}
\NormalTok{    merged }\OperatorTok{=}\NormalTok{ add\_derived\_metrics(merged)}
\NormalTok{    gdf }\OperatorTok{=}\NormalTok{ gpd.GeoDataFrame(merged, geometry}\OperatorTok{=}\StringTok{"geometry"}\NormalTok{, crs}\OperatorTok{=}\NormalTok{civ.crs)}
\NormalTok{    gdf.to\_file(config.CIVIC\_COMBINED, driver}\OperatorTok{=}\StringTok{"GeoJSON"}\NormalTok{)}
    \BuiltInTok{print}\NormalTok{(}\SpecialStringTok{f"Wrote }\SpecialCharTok{\{}\BuiltInTok{len}\NormalTok{(gdf)}\SpecialCharTok{\}}\SpecialStringTok{ combined civic{-}association rows → }\SpecialCharTok{\{}\NormalTok{config}\SpecialCharTok{.}\NormalTok{CIVIC\_COMBINED}\SpecialCharTok{\}}\SpecialStringTok{"}\NormalTok{)}


\ControlFlowTok{if} \VariableTok{\_\_name\_\_} \OperatorTok{==} \StringTok{"\_\_main\_\_"}\NormalTok{:}
\NormalTok{    main()}
\end{Highlighting}
\end{Shaded}

\section{Acquire parcels}\label{acquire-parcels}

Page through the Arlington MHUD FeatureServer and reduce parcel polygons
to unit-bearing centroids.

\textbf{\texttt{pipeline/acquire\_parcels.py}}

\begin{Shaded}
\begin{Highlighting}[]
\CommentTok{"""Acquire Arlington MHUD parcel polygons and reduce them to unit{-}bearing centroids."""}
\ImportTok{from}\NormalTok{ \_\_future\_\_ }\ImportTok{import}\NormalTok{ annotations}

\ImportTok{import}\NormalTok{ json}
\ImportTok{import}\NormalTok{ urllib.parse}
\ImportTok{import}\NormalTok{ urllib.request}

\ImportTok{import}\NormalTok{ geopandas }\ImportTok{as}\NormalTok{ gpd}
\ImportTok{import}\NormalTok{ pandas }\ImportTok{as}\NormalTok{ pd}

\ImportTok{from}\NormalTok{ pipeline }\ImportTok{import}\NormalTok{ config}

\NormalTok{LAYER }\OperatorTok{=}\NormalTok{ (}
    \StringTok{"https://arlgis.arlingtonva.us/arcgis/rest/services/Open\_Data/"}
    \StringTok{"od\_MHUD\_Polygons/FeatureServer/0/query"}
\NormalTok{)}
\NormalTok{PAGE }\OperatorTok{=} \DecValTok{2000}


\KeywordTok{def}\NormalTok{ \_fetch\_page(offset: }\BuiltInTok{int}\NormalTok{) }\OperatorTok{{-}\textgreater{}}\NormalTok{ gpd.GeoDataFrame }\OperatorTok{|} \VariableTok{None}\NormalTok{:}
\NormalTok{    params }\OperatorTok{=}\NormalTok{ \{}
        \StringTok{"where"}\NormalTok{: }\StringTok{"1=1"}\NormalTok{,}
        \StringTok{"outFields"}\NormalTok{: }\StringTok{"Total\_Units,Unit\_Type"}\NormalTok{,}
        \StringTok{"outSR"}\NormalTok{: }\StringTok{"4326"}\NormalTok{,}
        \StringTok{"resultOffset"}\NormalTok{: }\BuiltInTok{str}\NormalTok{(offset),}
        \StringTok{"resultRecordCount"}\NormalTok{: }\BuiltInTok{str}\NormalTok{(PAGE),}
        \StringTok{"f"}\NormalTok{: }\StringTok{"geojson"}\NormalTok{,}
\NormalTok{    \}}
\NormalTok{    url }\OperatorTok{=}\NormalTok{ LAYER }\OperatorTok{+} \StringTok{"?"} \OperatorTok{+}\NormalTok{ urllib.parse.urlencode(params)}
    \ControlFlowTok{with}\NormalTok{ urllib.request.urlopen(url, timeout}\OperatorTok{=}\DecValTok{60}\NormalTok{) }\ImportTok{as}\NormalTok{ resp:  }\CommentTok{\# noqa: S310}
\NormalTok{        data }\OperatorTok{=}\NormalTok{ json.load(resp)}
\NormalTok{    feats }\OperatorTok{=}\NormalTok{ data.get(}\StringTok{"features"}\NormalTok{, [])}
    \ControlFlowTok{if} \KeywordTok{not}\NormalTok{ feats:}
        \ControlFlowTok{return} \VariableTok{None}
    \ControlFlowTok{return}\NormalTok{ gpd.GeoDataFrame.from\_features(feats, crs}\OperatorTok{=}\StringTok{"EPSG:4326"}\NormalTok{)}


\KeywordTok{def}\NormalTok{ fetch\_parcels() }\OperatorTok{{-}\textgreater{}}\NormalTok{ gpd.GeoDataFrame:}
    \CommentTok{"""Page through the MHUD FeatureServer and return all parcel polygons."""}
\NormalTok{    parts: }\BuiltInTok{list}\NormalTok{[gpd.GeoDataFrame] }\OperatorTok{=}\NormalTok{ []}
\NormalTok{    offset }\OperatorTok{=} \DecValTok{0}
    \ControlFlowTok{while} \VariableTok{True}\NormalTok{:}
\NormalTok{        page }\OperatorTok{=}\NormalTok{ \_fetch\_page(offset)}
        \ControlFlowTok{if}\NormalTok{ page }\KeywordTok{is} \VariableTok{None} \KeywordTok{or} \BuiltInTok{len}\NormalTok{(page) }\OperatorTok{==} \DecValTok{0}\NormalTok{:}
            \ControlFlowTok{break}
\NormalTok{        parts.append(page)}
        \ControlFlowTok{if} \BuiltInTok{len}\NormalTok{(page) }\OperatorTok{\textless{}}\NormalTok{ PAGE:}
            \ControlFlowTok{break}
\NormalTok{        offset }\OperatorTok{+=}\NormalTok{ PAGE}
\NormalTok{    gdf }\OperatorTok{=}\NormalTok{ gpd.GeoDataFrame(pd.concat(parts, ignore\_index}\OperatorTok{=}\VariableTok{True}\NormalTok{), crs}\OperatorTok{=}\StringTok{"EPSG:4326"}\NormalTok{)}
    \ControlFlowTok{return}\NormalTok{ gdf}


\KeywordTok{def}\NormalTok{ to\_centroids(polys: gpd.GeoDataFrame) }\OperatorTok{{-}\textgreater{}}\NormalTok{ gpd.GeoDataFrame:}
    \CommentTok{"""Reduce parcel polygons to centroids, keeping Total\_Units and Unit\_Type."""}
\NormalTok{    cent }\OperatorTok{=}\NormalTok{ polys.copy()}
    \CommentTok{\# representative\_point() is guaranteed inside the polygon (unlike centroid)}
\NormalTok{    cent[}\StringTok{"geometry"}\NormalTok{] }\OperatorTok{=}\NormalTok{ cent.geometry.representative\_point()}
    \ControlFlowTok{return}\NormalTok{ cent[[}\StringTok{"Total\_Units"}\NormalTok{, }\StringTok{"Unit\_Type"}\NormalTok{, }\StringTok{"geometry"}\NormalTok{]]}


\KeywordTok{def}\NormalTok{ main() }\OperatorTok{{-}\textgreater{}} \VariableTok{None}\NormalTok{:}
\NormalTok{    polys }\OperatorTok{=}\NormalTok{ fetch\_parcels()}
\NormalTok{    pts }\OperatorTok{=}\NormalTok{ to\_centroids(polys)}
\NormalTok{    pts.to\_file(config.PARCELS, driver}\OperatorTok{=}\StringTok{"GeoJSON"}\NormalTok{)}
    \BuiltInTok{print}\NormalTok{(}\SpecialStringTok{f"Wrote }\SpecialCharTok{\{}\BuiltInTok{len}\NormalTok{(pts)}\SpecialCharTok{\}}\SpecialStringTok{ parcel centroids → }\SpecialCharTok{\{}\NormalTok{config}\SpecialCharTok{.}\NormalTok{PARCELS}\SpecialCharTok{\}}\SpecialStringTok{"}\NormalTok{)}


\ControlFlowTok{if} \VariableTok{\_\_name\_\_} \OperatorTok{==} \StringTok{"\_\_main\_\_"}\NormalTok{:}
\NormalTok{    main()}
\end{Highlighting}
\end{Shaded}

\section{Redistribute income through
parcels}\label{redistribute-income-through-parcels}

Explode parcels into per-unit points and redistribute block-group income
through them (the dasymetric second approach).

\textbf{\texttt{pipeline/redistribute\_income\_parcels.py}}

\begin{Shaded}
\begin{Highlighting}[]
\CommentTok{"""Redistribute income through unit{-}weighted parcels (dasymetric), derive mean income."""}
\ImportTok{from}\NormalTok{ \_\_future\_\_ }\ImportTok{import}\NormalTok{ annotations}

\ImportTok{from}\NormalTok{ pathlib }\ImportTok{import}\NormalTok{ Path}

\ImportTok{import}\NormalTok{ geopandas }\ImportTok{as}\NormalTok{ gpd}
\ImportTok{import}\NormalTok{ pandas }\ImportTok{as}\NormalTok{ pd}
\ImportTok{from}\NormalTok{ sdc\_redistribute }\ImportTok{import}\NormalTok{ redistribute\_parcels}

\ImportTok{from}\NormalTok{ pipeline }\ImportTok{import}\NormalTok{ config}


\KeywordTok{def}\NormalTok{ explode\_by\_units(parcels: gpd.GeoDataFrame) }\OperatorTok{{-}\textgreater{}}\NormalTok{ gpd.GeoDataFrame:}
    \CommentTok{"""Replicate each parcel into Total\_Units identical points; drop units \textless{}= 0."""}
\NormalTok{    p }\OperatorTok{=}\NormalTok{ parcels[parcels[}\StringTok{"Total\_Units"}\NormalTok{] }\OperatorTok{\textgreater{}} \DecValTok{0}\NormalTok{].copy()}
\NormalTok{    reps }\OperatorTok{=}\NormalTok{ p[}\StringTok{"Total\_Units"}\NormalTok{].astype(}\BuiltInTok{int}\NormalTok{).to\_numpy()}
\NormalTok{    out }\OperatorTok{=}\NormalTok{ p.loc[p.index.repeat(reps)].reset\_index(drop}\OperatorTok{=}\VariableTok{True}\NormalTok{)}
    \ControlFlowTok{return}\NormalTok{ out}


\KeywordTok{def}\NormalTok{ redistribute\_income\_parcels(}
\NormalTok{    counts\_long: pd.DataFrame,}
\NormalTok{    parcels: gpd.GeoDataFrame,}
\NormalTok{    source\_geo: Path,}
\NormalTok{    target\_geo: Path,}
\NormalTok{) }\OperatorTok{{-}\textgreater{}}\NormalTok{ pd.DataFrame:}
    \CommentTok{"""Return per{-}civic{-}association agg\_income, households, derived mean\_income."""}
\NormalTok{    points }\OperatorTok{=}\NormalTok{ explode\_by\_units(parcels)}
\NormalTok{    result }\OperatorTok{=}\NormalTok{ redistribute\_parcels(}
\NormalTok{        source\_df}\OperatorTok{=}\NormalTok{counts\_long,}
\NormalTok{        parcel\_centroids}\OperatorTok{=}\NormalTok{points,}
\NormalTok{        source\_geo}\OperatorTok{=}\NormalTok{source\_geo,}
\NormalTok{        target\_geos}\OperatorTok{=}\NormalTok{\{}\StringTok{"civic\_association"}\NormalTok{: target\_geo\},}
\NormalTok{        count\_cols}\OperatorTok{=}\NormalTok{[}\StringTok{"agg\_income"}\NormalTok{, }\StringTok{"households"}\NormalTok{],}
\NormalTok{        source\_id}\OperatorTok{=}\StringTok{"geoid"}\NormalTok{,}
\NormalTok{    )}
\NormalTok{    wide }\OperatorTok{=}\NormalTok{ result.pivot\_table(}
\NormalTok{        index}\OperatorTok{=}\StringTok{"geoid"}\NormalTok{, columns}\OperatorTok{=}\StringTok{"measure"}\NormalTok{, values}\OperatorTok{=}\StringTok{"value"}\NormalTok{, aggfunc}\OperatorTok{=}\StringTok{"first"}
\NormalTok{    ).reset\_index()}
\NormalTok{    wide.columns.name }\OperatorTok{=} \VariableTok{None}
\NormalTok{    wide }\OperatorTok{=}\NormalTok{ wide.rename(}
\NormalTok{        columns}\OperatorTok{=}\NormalTok{\{}\StringTok{"agg\_income\_parcels"}\NormalTok{: }\StringTok{"agg\_income"}\NormalTok{, }\StringTok{"households\_parcels"}\NormalTok{: }\StringTok{"households"}\NormalTok{\}}
\NormalTok{    )}
\NormalTok{    wide[}\StringTok{"mean\_income"}\NormalTok{] }\OperatorTok{=}\NormalTok{ wide[}\StringTok{"agg\_income"}\NormalTok{] }\OperatorTok{/}\NormalTok{ wide[}\StringTok{"households"}\NormalTok{].where(}
\NormalTok{        wide[}\StringTok{"households"}\NormalTok{] }\OperatorTok{\textgreater{}} \DecValTok{0}
\NormalTok{    )}
    \ControlFlowTok{return}\NormalTok{ wide[[}\StringTok{"geoid"}\NormalTok{, }\StringTok{"agg\_income"}\NormalTok{, }\StringTok{"households"}\NormalTok{, }\StringTok{"mean\_income"}\NormalTok{]]}


\KeywordTok{def}\NormalTok{ main() }\OperatorTok{{-}\textgreater{}} \VariableTok{None}\NormalTok{:}
\NormalTok{    parcels }\OperatorTok{=}\NormalTok{ gpd.read\_file(config.PARCELS)}
\NormalTok{    bg }\OperatorTok{=}\NormalTok{ gpd.read\_file(config.ACS\_COUNTS).rename(columns}\OperatorTok{=}\NormalTok{\{}\StringTok{"GEOID"}\NormalTok{: }\StringTok{"geoid"}\NormalTok{\})}
    \BuiltInTok{long} \OperatorTok{=}\NormalTok{ bg[[}\StringTok{"geoid"}\NormalTok{, }\StringTok{"agg\_income"}\NormalTok{, }\StringTok{"households"}\NormalTok{]].melt(}
\NormalTok{        id\_vars}\OperatorTok{=}\StringTok{"geoid"}\NormalTok{, var\_name}\OperatorTok{=}\StringTok{"measure"}\NormalTok{, value\_name}\OperatorTok{=}\StringTok{"value"}
\NormalTok{    )}
    \BuiltInTok{long}\NormalTok{[}\StringTok{"year"}\NormalTok{] }\OperatorTok{=}\NormalTok{ config.ACS\_YEAR}

\NormalTok{    src\_path }\OperatorTok{=}\NormalTok{ config.DATA }\OperatorTok{/} \StringTok{"\_acs\_src\_parcels.geojson"}
\NormalTok{    bg[[}\StringTok{"geoid"}\NormalTok{, }\StringTok{"geometry"}\NormalTok{]].to\_file(src\_path, driver}\OperatorTok{=}\StringTok{"GeoJSON"}\NormalTok{)}
\NormalTok{    out }\OperatorTok{=}\NormalTok{ redistribute\_income\_parcels(}\BuiltInTok{long}\NormalTok{, parcels, src\_path, config.CIVIC\_ASSOC)}
\NormalTok{    src\_path.unlink(missing\_ok}\OperatorTok{=}\VariableTok{True}\NormalTok{)}

\NormalTok{    out.to\_csv(config.CIVIC\_INCOME\_PARCELS, index}\OperatorTok{=}\VariableTok{False}\NormalTok{)}
    \BuiltInTok{print}\NormalTok{(}\SpecialStringTok{f"Wrote }\SpecialCharTok{\{}\BuiltInTok{len}\NormalTok{(out)}\SpecialCharTok{\}}\SpecialStringTok{ parcel{-}weighted income rows → }\SpecialCharTok{\{}\NormalTok{config}\SpecialCharTok{.}\NormalTok{CIVIC\_INCOME\_PARCELS}\SpecialCharTok{\}}\SpecialStringTok{"}\NormalTok{)}


\ControlFlowTok{if} \VariableTok{\_\_name\_\_} \OperatorTok{==} \StringTok{"\_\_main\_\_"}\NormalTok{:}
\NormalTok{    main()}
\end{Highlighting}
\end{Shaded}

\section{Compare methods}\label{compare-methods}

Join the area-weighted and parcel-weighted income and quantify the
difference per civic association.

\textbf{\texttt{pipeline/compare\_methods.py}}

\begin{Shaded}
\begin{Highlighting}[]
\CommentTok{"""Compare area{-}weighted vs parcel{-}weighted income per civic association."""}
\ImportTok{from}\NormalTok{ \_\_future\_\_ }\ImportTok{import}\NormalTok{ annotations}

\ImportTok{import}\NormalTok{ geopandas }\ImportTok{as}\NormalTok{ gpd}
\ImportTok{import}\NormalTok{ pandas }\ImportTok{as}\NormalTok{ pd}

\ImportTok{from}\NormalTok{ pipeline }\ImportTok{import}\NormalTok{ config}


\KeywordTok{def}\NormalTok{ compare\_income(area: pd.DataFrame, parcel: pd.DataFrame) }\OperatorTok{{-}\textgreater{}}\NormalTok{ pd.DataFrame:}
    \CommentTok{"""Join the two methods\textquotesingle{} mean income; add diff (parcel {-} area) and pct\_diff."""}
\NormalTok{    a }\OperatorTok{=}\NormalTok{ area[[}\StringTok{"geoid"}\NormalTok{, }\StringTok{"mean\_income"}\NormalTok{]].rename(columns}\OperatorTok{=}\NormalTok{\{}\StringTok{"mean\_income"}\NormalTok{: }\StringTok{"mean\_income\_area"}\NormalTok{\})}
\NormalTok{    p }\OperatorTok{=}\NormalTok{ parcel[[}\StringTok{"geoid"}\NormalTok{, }\StringTok{"mean\_income"}\NormalTok{]].rename(columns}\OperatorTok{=}\NormalTok{\{}\StringTok{"mean\_income"}\NormalTok{: }\StringTok{"mean\_income\_parcel"}\NormalTok{\})}
\NormalTok{    out }\OperatorTok{=}\NormalTok{ a.merge(p, on}\OperatorTok{=}\StringTok{"geoid"}\NormalTok{, how}\OperatorTok{=}\StringTok{"outer"}\NormalTok{)}
\NormalTok{    out[}\StringTok{"diff"}\NormalTok{] }\OperatorTok{=}\NormalTok{ out[}\StringTok{"mean\_income\_parcel"}\NormalTok{] }\OperatorTok{{-}}\NormalTok{ out[}\StringTok{"mean\_income\_area"}\NormalTok{]}
\NormalTok{    out[}\StringTok{"pct\_diff"}\NormalTok{] }\OperatorTok{=} \DecValTok{100} \OperatorTok{*}\NormalTok{ out[}\StringTok{"diff"}\NormalTok{] }\OperatorTok{/}\NormalTok{ out[}\StringTok{"mean\_income\_area"}\NormalTok{].where(}
\NormalTok{        out[}\StringTok{"mean\_income\_area"}\NormalTok{] }\OperatorTok{\textgreater{}} \DecValTok{0}
\NormalTok{    )}
    \ControlFlowTok{return}\NormalTok{ out}


\KeywordTok{def}\NormalTok{ main() }\OperatorTok{{-}\textgreater{}} \VariableTok{None}\NormalTok{:}
\NormalTok{    civ }\OperatorTok{=}\NormalTok{ gpd.read\_file(config.CIVIC\_ASSOC)[[}\StringTok{"geoid"}\NormalTok{, }\StringTok{"region\_name"}\NormalTok{, }\StringTok{"geometry"}\NormalTok{]]}
\NormalTok{    area }\OperatorTok{=}\NormalTok{ pd.read\_csv(config.CIVIC\_INCOME, dtype}\OperatorTok{=}\NormalTok{\{}\StringTok{"geoid"}\NormalTok{: }\BuiltInTok{str}\NormalTok{\})}
\NormalTok{    parcel }\OperatorTok{=}\NormalTok{ pd.read\_csv(config.CIVIC\_INCOME\_PARCELS, dtype}\OperatorTok{=}\NormalTok{\{}\StringTok{"geoid"}\NormalTok{: }\BuiltInTok{str}\NormalTok{\})}
    \BuiltInTok{cmp} \OperatorTok{=}\NormalTok{ compare\_income(area, parcel)}
\NormalTok{    gdf }\OperatorTok{=}\NormalTok{ civ.merge(}\BuiltInTok{cmp}\NormalTok{, on}\OperatorTok{=}\StringTok{"geoid"}\NormalTok{, how}\OperatorTok{=}\StringTok{"left"}\NormalTok{)}
\NormalTok{    gdf }\OperatorTok{=}\NormalTok{ gpd.GeoDataFrame(gdf, geometry}\OperatorTok{=}\StringTok{"geometry"}\NormalTok{, crs}\OperatorTok{=}\NormalTok{civ.crs)}

    \CommentTok{\# Honest accounting: report parcel{-}method household total vs ACS source.}
\NormalTok{    acs }\OperatorTok{=}\NormalTok{ gpd.read\_file(config.ACS\_COUNTS)}
\NormalTok{    src\_hh }\OperatorTok{=} \BuiltInTok{float}\NormalTok{(acs[}\StringTok{"households"}\NormalTok{].}\BuiltInTok{sum}\NormalTok{())}
\NormalTok{    parcel\_hh }\OperatorTok{=} \BuiltInTok{float}\NormalTok{(pd.read\_csv(config.CIVIC\_INCOME\_PARCELS)[}\StringTok{"households"}\NormalTok{].}\BuiltInTok{sum}\NormalTok{())}
\NormalTok{    pct }\OperatorTok{=} \DecValTok{100} \OperatorTok{*}\NormalTok{ (parcel\_hh }\OperatorTok{{-}}\NormalTok{ src\_hh) }\OperatorTok{/}\NormalTok{ src\_hh }\ControlFlowTok{if}\NormalTok{ src\_hh }\ControlFlowTok{else} \FloatTok{0.0}
    \BuiltInTok{print}\NormalTok{(}\SpecialStringTok{f"Household totals, ACS source: }\SpecialCharTok{\{}\NormalTok{src\_hh}\SpecialCharTok{:,.0f\}}\SpecialStringTok{; parcel method: }\SpecialCharTok{\{}\NormalTok{parcel\_hh}\SpecialCharTok{:,.0f\}}\SpecialStringTok{ (}\SpecialCharTok{\{}\NormalTok{pct}\SpecialCharTok{:+.1f\}}\SpecialStringTok{\%)"}\NormalTok{)}

\NormalTok{    gdf.to\_file(config.CIVIC\_INCOME\_COMPARISON, driver}\OperatorTok{=}\StringTok{"GeoJSON"}\NormalTok{)}
    \BuiltInTok{print}\NormalTok{(}\SpecialStringTok{f"Wrote }\SpecialCharTok{\{}\BuiltInTok{len}\NormalTok{(gdf)}\SpecialCharTok{\}}\SpecialStringTok{ comparison rows → }\SpecialCharTok{\{}\NormalTok{config}\SpecialCharTok{.}\NormalTok{CIVIC\_INCOME\_COMPARISON}\SpecialCharTok{\}}\SpecialStringTok{"}\NormalTok{)}


\ControlFlowTok{if} \VariableTok{\_\_name\_\_} \OperatorTok{==} \StringTok{"\_\_main\_\_"}\NormalTok{:}
\NormalTok{    main()}
\end{Highlighting}
\end{Shaded}

\section{Validate}\label{validate}

Confirm household totals are preserved by the redistribution and that no
speeds or incomes are negative.

\textbf{\texttt{pipeline/validate.py}}

\begin{Shaded}
\begin{Highlighting}[]
\CommentTok{"""Integrity checks on redistributed outputs."""}
\ImportTok{from}\NormalTok{ \_\_future\_\_ }\ImportTok{import}\NormalTok{ annotations}

\ImportTok{import}\NormalTok{ pandas }\ImportTok{as}\NormalTok{ pd}


\KeywordTok{def}\NormalTok{ check\_total\_preserved(source\_total: }\BuiltInTok{float}\NormalTok{, target\_total: }\BuiltInTok{float}\NormalTok{, tol: }\BuiltInTok{float} \OperatorTok{=} \FloatTok{0.01}\NormalTok{) }\OperatorTok{{-}\textgreater{}} \VariableTok{None}\NormalTok{:}
    \CommentTok{"""Assert target total matches source total within relative tolerance \textasciigrave{}\textasciigrave{}tol\textasciigrave{}\textasciigrave{}."""}
    \ControlFlowTok{if}\NormalTok{ source\_total }\OperatorTok{==} \DecValTok{0}\NormalTok{:}
        \ControlFlowTok{return}
\NormalTok{    rel }\OperatorTok{=} \BuiltInTok{abs}\NormalTok{(target\_total }\OperatorTok{{-}}\NormalTok{ source\_total) }\OperatorTok{/} \BuiltInTok{abs}\NormalTok{(source\_total)}
    \ControlFlowTok{assert}\NormalTok{ rel }\OperatorTok{\textless{}=}\NormalTok{ tol, }\SpecialStringTok{f"total drift }\SpecialCharTok{\{}\NormalTok{rel}\SpecialCharTok{:.3\%\}}\SpecialStringTok{ exceeds }\SpecialCharTok{\{}\NormalTok{tol}\SpecialCharTok{:.3\%\}}\SpecialStringTok{"}


\KeywordTok{def}\NormalTok{ check\_ranges(df: pd.DataFrame) }\OperatorTok{{-}\textgreater{}} \VariableTok{None}\NormalTok{:}
    \CommentTok{"""Assert speed/income columns are non{-}negative where present."""}
    \ControlFlowTok{for}\NormalTok{ col }\KeywordTok{in}\NormalTok{ (}\StringTok{"download\_mbps"}\NormalTok{, }\StringTok{"upload\_mbps"}\NormalTok{, }\StringTok{"mean\_income"}\NormalTok{, }\StringTok{"households"}\NormalTok{):}
        \ControlFlowTok{if}\NormalTok{ col }\KeywordTok{in}\NormalTok{ df.columns:}
\NormalTok{            vals }\OperatorTok{=}\NormalTok{ df[col].dropna()}
            \ControlFlowTok{assert}\NormalTok{ (vals }\OperatorTok{\textgreater{}=} \DecValTok{0}\NormalTok{).}\BuiltInTok{all}\NormalTok{(), }\SpecialStringTok{f"}\SpecialCharTok{\{}\NormalTok{col}\SpecialCharTok{\}}\SpecialStringTok{ has negative values"}


\KeywordTok{def}\NormalTok{ main() }\OperatorTok{{-}\textgreater{}} \VariableTok{None}\NormalTok{:}
    \ImportTok{import}\NormalTok{ geopandas }\ImportTok{as}\NormalTok{ gpd}
    \ImportTok{from}\NormalTok{ pipeline }\ImportTok{import}\NormalTok{ config}

\NormalTok{    combined }\OperatorTok{=}\NormalTok{ gpd.read\_file(config.CIVIC\_COMBINED)}
\NormalTok{    acs }\OperatorTok{=}\NormalTok{ gpd.read\_file(config.ACS\_COUNTS)}
\NormalTok{    check\_total\_preserved(}
\NormalTok{        acs[}\StringTok{"households"}\NormalTok{].}\BuiltInTok{sum}\NormalTok{(), combined[}\StringTok{"households"}\NormalTok{].}\BuiltInTok{sum}\NormalTok{(), tol}\OperatorTok{=}\FloatTok{0.02}
\NormalTok{    )}
\NormalTok{    check\_ranges(combined)}
    \BuiltInTok{print}\NormalTok{(}\StringTok{"Validation passed."}\NormalTok{)}


\ControlFlowTok{if} \VariableTok{\_\_name\_\_} \OperatorTok{==} \StringTok{"\_\_main\_\_"}\NormalTok{:}
\NormalTok{    main()}
\end{Highlighting}
\end{Shaded}

\section{Run the full pipeline}\label{run-the-full-pipeline}

Wire every stage into one command:
\texttt{uv\ run\ python\ -m\ pipeline.run}.

\textbf{\texttt{pipeline/run.py}}

\begin{Shaded}
\begin{Highlighting}[]
\CommentTok{"""Run the full pipeline end to end (acquisition → outputs → validation)."""}
\ImportTok{from}\NormalTok{ \_\_future\_\_ }\ImportTok{import}\NormalTok{ annotations}

\ImportTok{from}\NormalTok{ pipeline }\ImportTok{import}\NormalTok{ (}
\NormalTok{    acquire\_acs,}
\NormalTok{    acquire\_geographies,}
\NormalTok{    acquire\_ookla,}
\NormalTok{    acquire\_parcels,}
\NormalTok{    combine,}
\NormalTok{    compare\_methods,}
\NormalTok{    redistribute\_broadband,}
\NormalTok{    redistribute\_income,}
\NormalTok{    redistribute\_income\_parcels,}
\NormalTok{    validate,}
\NormalTok{)}


\KeywordTok{def}\NormalTok{ main() }\OperatorTok{{-}\textgreater{}} \VariableTok{None}\NormalTok{:}
\NormalTok{    acquire\_geographies.main()}
\NormalTok{    acquire\_acs.main()}
\NormalTok{    acquire\_ookla.main()}
\NormalTok{    acquire\_parcels.main()}
\NormalTok{    redistribute\_income.main()}
\NormalTok{    redistribute\_broadband.main()}
\NormalTok{    combine.main()}
\NormalTok{    redistribute\_income\_parcels.main()}
\NormalTok{    compare\_methods.main()}
\NormalTok{    validate.main()}
    \BuiltInTok{print}\NormalTok{(}\StringTok{"Pipeline complete."}\NormalTok{)}


\ControlFlowTok{if} \VariableTok{\_\_name\_\_} \OperatorTok{==} \StringTok{"\_\_main\_\_"}\NormalTok{:}
\NormalTok{    main()}
\end{Highlighting}
\end{Shaded}

\bookmarksetup{startatroot}

\chapter*{References}\label{references}
\addcontentsline{toc}{chapter}{References}

\markboth{References}{References}

\phantomsection\label{refs}
\begin{CSLReferences}{0}{0}
\bibitem[\citeproctext]{ref-ref1}
\CSLLeftMargin{{[}1{]} }%
\CSLRightInline{{``Stacks.cdc.gov.''} Accessed: Jun. 04, 2026.
{[}Online{]}. Available:
\url{https://stacks.cdc.gov/view/cdc/59750/cdc_59750_DS1.pdf}}

\bibitem[\citeproctext]{ref-ref4}
\CSLLeftMargin{{[}2{]} }%
\CSLRightInline{{``Github.com.''} Accessed: Jun. 04, 2026. {[}Online{]}.
Available: \url{https://github.com/CDCgov/EPHTracking-Subcounty}}

\bibitem[\citeproctext]{ref-ref32}
\CSLLeftMargin{{[}3{]} }%
\CSLRightInline{{``results4america.org.''} Accessed: Jun. 04, 2026.
{[}Online{]}. Available:
\url{https://results4america.org/wp-content/uploads/2021/06/Deloitte-WWC-Data-Gap-Report_vFinal-063021.pdf}}

\bibitem[\citeproctext]{ref-ref45}
\CSLLeftMargin{{[}4{]} }%
\CSLRightInline{{``Iecam.illinois.edu.''} Accessed: Jun. 04, 2026.
{[}Online{]}. Available:
\url{https://iecam.illinois.edu/browse/subcounty-data-cautions-and-recommendations}}

\bibitem[\citeproctext]{ref-ref50}
\CSLLeftMargin{{[}5{]} }%
\CSLRightInline{{``Atlas.co.''} Accessed: Jun. 04, 2026. {[}Online{]}.
Available:
\url{https://atlas.co/blog/modifiable-areal-unit-problem-maup/}}

\bibitem[\citeproctext]{ref-ref79}
\CSLLeftMargin{{[}6{]} }%
\CSLRightInline{{``Gisgeography.com.''} Accessed: Jun. 04, 2026.
{[}Online{]}. Available:
\url{https://gisgeography.com/maup-modifiable-areal-unit-problem/}}

\bibitem[\citeproctext]{ref-ref86}
\CSLLeftMargin{{[}7{]} }%
\CSLRightInline{{``Speedtest by ookla global fixed and mobile network
performance ...''} Accessed: Jun. 04, 2026. {[}Online{]}. Available:
\url{https://github.com/teamookla/ookla-open-data}}

\bibitem[\citeproctext]{ref-ref88}
\CSLLeftMargin{{[}8{]} }%
\CSLRightInline{{``Announcing ookla open datasets.''} Accessed: Jun. 04,
2026. {[}Online{]}. Available:
\url{https://www.ookla.com/articles/announcing-ookla-open-datasets}}

\bibitem[\citeproctext]{ref-ref89}
\CSLLeftMargin{{[}9{]} }%
\CSLRightInline{{``American community survey 5-year data (2009-2023).''}
Accessed: Jun. 04, 2026. {[}Online{]}. Available:
\url{https://www.census.gov/data/developers/data-sets/acs-5year.html}}

\bibitem[\citeproctext]{ref-ref142}
\CSLLeftMargin{{[}10{]} }%
\CSLRightInline{{``American community survey data - u.s. Census
bureau.''} Accessed: Jun. 04, 2026. {[}Online{]}. Available:
\url{https://www.census.gov/programs-surveys/acs/data.html}}

\bibitem[\citeproctext]{ref-ref91}
\CSLLeftMargin{{[}11{]} }%
\CSLRightInline{{``Arlington county civic federation: home.''} Accessed:
Jun. 04, 2026. {[}Online{]}. Available: \url{https://www.civfed.org}}

\bibitem[\citeproctext]{ref-ref145}
\CSLLeftMargin{{[}12{]} }%
\CSLRightInline{{``Civic \& citizen associations - arlington county.''}
Accessed: Jun. 04, 2026. {[}Online{]}. Available:
\url{https://www.arlingtonva.us/Government/Topics/Community/Civic-and-Citizen-Association}}

\bibitem[\citeproctext]{ref-ref146}
\CSLLeftMargin{{[}13{]} }%
\CSLRightInline{{``Civic association polygons \textbar{} arlington
county, virginia. GIS open data.''} Accessed: Jun. 04, 2026.
{[}Online{]}. Available:
\url{https://gisdata-arlgis.opendata.arcgis.com/datasets/ArlGIS::civic-association-polygons/explore?location=38.880938\%2C-77.109550\%2C13.55}}

\bibitem[\citeproctext]{ref-ref99}
\CSLLeftMargin{{[}14{]} }%
\CSLRightInline{{``What is areal interpolation?---ArcGIS pro \textbar{}
documentation.''} Accessed: Jun. 04, 2026. {[}Online{]}. Available:
\url{https://pro.arcgis.com/en/pro-app/latest/help/analysis/geostatistical-analyst/what-is-areal-interpolation.htm}}

\bibitem[\citeproctext]{ref-ref135}
\CSLLeftMargin{{[}15{]} }%
\CSLRightInline{{``Areal interpolation \textbar{} vincent thorne.''}
Accessed: Jun. 04, 2026. {[}Online{]}. Available:
\url{https://vinceth.net/2021/06/18/areal-interpolation.html}}

\bibitem[\citeproctext]{ref-ref139}
\CSLLeftMargin{{[}16{]} }%
\CSLRightInline{{``Spatial interpolation using areal features: A review
of methods and ...''} Accessed: Jun. 04, 2026. {[}Online{]}. Available:
\url{https://compass.onlinelibrary.wiley.com/doi/10.1111/gec3.12465}}

\bibitem[\citeproctext]{ref-ref159}
\CSLLeftMargin{{[}17{]} }%
\CSLRightInline{{``Arlington county, virginia - u.s. Census bureau
QuickFacts.''} Accessed: Jun. 04, 2026. {[}Online{]}. Available:
\url{https://www.census.gov/quickfacts/fact/table/arlingtoncountyvirginia/PST045224}}

\bibitem[\citeproctext]{ref-ref95}
\CSLLeftMargin{{[}18{]} }%
\CSLRightInline{{``Mapping the digital divide: What predicts internet
access across ...''} Accessed: Jun. 04, 2026. {[}Online{]}. Available:
\url{https://pmc.ncbi.nlm.nih.gov/articles/PMC10961931/}}

\bibitem[\citeproctext]{ref-ref96}
\CSLLeftMargin{{[}19{]} }%
\CSLRightInline{{``Bridging the digital divide: A path to universal
broadband access.''} Accessed: Jun. 04, 2026. {[}Online{]}. Available:
\url{https://broadbandbreakfast.com/bridging-the-digital-divide-a-path-to-universal-broadband-access/}}

\bibitem[\citeproctext]{ref-ref112}
\CSLLeftMargin{{[}20{]} }%
\CSLRightInline{{``How policymakers can help bridge the digital divide
in 2021 - AAF.''} Accessed: Jun. 04, 2026. {[}Online{]}. Available:
\url{https://www.americanactionforum.org/insight/how-policymakers-can-help-bridge-the-digital-divide-in-2021/}}

\bibitem[\citeproctext]{ref-ref100}
\CSLLeftMargin{{[}21{]} }%
\CSLRightInline{{``Areal interpolation \textbar{} by ITC, university of
twente.''} Accessed: Jun. 04, 2026. {[}Online{]}. Available:
\url{https://gistbok-ltb.ucgis.org/27/concept/7992}}

\bibitem[\citeproctext]{ref-ref137}
\CSLLeftMargin{{[}22{]} }%
\CSLRightInline{{``Speedtest® methodology \textbar{} ookla®.''}
Accessed: Jun. 04, 2026. {[}Online{]}. Available:
\url{https://www.ookla.com/resources/guides/speedtest-methodology}}

\bibitem[\citeproctext]{ref-ref156}
\CSLLeftMargin{{[}23{]} }%
\CSLRightInline{{``American community survey (ACS) - arlington
county.''} Accessed: Jun. 04, 2026. {[}Online{]}. Available:
\url{https://www.arlingtonva.us/Government/Projects/Data-Research/Demographics/American-Community-Survey}}

\end{CSLReferences}




\end{document}
